% Single-source research edition. Compile with pdfLaTeX twice, or latexmk -pdf.
\documentclass[11pt,letterpaper]{article}
\usepackage[T1]{fontenc}
\usepackage[utf8]{inputenc}
\usepackage{lmodern}
\usepackage[margin=1in,headheight=15pt]{geometry}
\usepackage{amsmath,amssymb,amsthm,mathtools}
\usepackage{microtype}
\usepackage{xcolor,listings,booktabs,array,longtable,enumitem,fancyhdr}
\usepackage{hyperref}
\definecolor{deepblue}{RGB}{28,58,96}
\definecolor{softgray}{RGB}{245,247,249}
\definecolor{codegray}{RGB}{45,52,60}
\hypersetup{colorlinks=true,linkcolor=deepblue,citecolor=deepblue,urlcolor=deepblue,
 pdftitle={Beukers--Delaygue Supercongruences, Apery Truncations, and the Level-54 Attractor, Complete Edition},
 pdfauthor={Research edition prepared for Keio},bookmarksnumbered=true,bookmarksopen=false}
\pagestyle{fancy}\fancyhf{}
\newcommand{\runningpaper}{Hypergeometric supercongruences}
\lhead{\small\itshape\runningpaper}
\rhead{\small\scshape Research edition}\cfoot{\thepage}
\setlength{\emergencystretch}{1.3em}
\setcounter{tocdepth}{1}
\newtheorem{theorem}{Theorem}[section]
\newtheorem{proposition}[theorem]{Proposition}
\newtheorem{corollary}[theorem]{Corollary}
\newtheorem{lemma}[theorem]{Lemma}
\newtheorem{conjecture}[theorem]{Conjecture}
\theoremstyle{definition}
\newtheorem{definition}[theorem]{Definition}
\theoremstyle{remark}
\newtheorem{remark}[theorem]{Remark}
\lstset{language=Python,basicstyle=\ttfamily\scriptsize,columns=fullflexible,
 keepspaces=true,breaklines=true,showstringspaces=false,frame=single,
 rulecolor=\color{black!25},keywordstyle=\color{blue!50!black},
 commentstyle=\color{black!60},aboveskip=8pt,belowskip=8pt}
\lstdefinestyle{python}{language=Python,basicstyle=\ttfamily\scriptsize,
 keywordstyle=\color{deepblue}\bfseries,commentstyle=\color{green!40!black},
 stringstyle=\color{purple!70!black},frame=single,rulecolor=\color{black!25},
 backgroundcolor=\color{softgray},breaklines=true,showstringspaces=false,
 columns=fullflexible,keepspaces=true}
\newcommand{\paperanchors}[1]{%
 \renewcommand{\theHsection}{#1.\arabic{section}}%
 \renewcommand{\theHsubsection}{\theHsection.\arabic{subsection}}%
 \renewcommand{\theHequation}{#1.\arabic{equation}}%
 \renewcommand{\theHtable}{#1.\arabic{table}}%
 \renewcommand{\theHtheorem}{\theHsection.\arabic{theorem}}}
\newcommand{\openpaper}[5]{%
 \clearpage\renewcommand{\runningpaper}{#3}
 \setcounter{section}{0}\setcounter{subsection}{0}
 \setcounter{equation}{0}\setcounter{theorem}{0}\setcounter{table}{0}
 \renewcommand{\thesection}{\arabic{section}}\paperanchors{#1}
 \phantomsection\addcontentsline{toc}{part}{#1. #2}\label{paper:#1}
 \thispagestyle{plain}
 {\sffamily\small\color{deepblue} PART #1\par}\vspace{9pt}
 {\raggedright\LARGE\bfseries #2\par}\vspace{5pt}
 #5{\small #4\par}\vspace{9pt}\hrule\vspace{10pt}}
\newcommand{\biblioanchor}{\phantomsection\addcontentsline{toc}{section}{References}}
\newcommand{\resultheading}[2]{\par\medskip\noindent\textbf{#1}\hfill{\small #2}\par\nobreak\noindent}

\begin{document}
\thispagestyle{plain}
\pdfbookmark[0]{Title and principal results}{edition}
{\sffamily\small\color{deepblue} COMPLETE RESEARCH EDITION\par}
\vspace{12pt}
{\raggedright\Huge\bfseries Beukers--Delaygue Supercongruences, Ap\'ery Truncations, and the Level-54 Attractor\par}
\vspace{12pt}
{\large Prepared for Keio\par}
\smallskip
{\small 12 November 2026\par}
\vspace{12pt}\hrule\vspace{14pt}
\noindent\textbf{Abstract.}
We study the precision of truncated hypergeometric periods under Frobenius.
For powers of the central binomial coefficients, finite reflection followed
by an integral lift of the digit operator proves the fixed-precision and
all-level statements of Beukers and Delaygue. An elliptic crystal gives a
two-step congruence for Ap\'ery truncations on smooth supersingular disks.
At the level-54 attractor of AESZ4 we construct a tetrahedral projector,
identify its two-dimensional image with the weight-four modular factor,
and interpolate the factorial blocks to recover the truncated period modulo $p^3$.
The resulting argument applies at every good prime, including primes for
which the modular coefficient is divisible by $p$.

\medskip
\noindent\textbf{Principal results and their hypotheses}
\smallskip

\resultheading{Reflection and all lengths}{Parts I--II, pp.~\pageref{paper:I}, \pageref{paper:II}}
For odd $p$, $d\ge2$, $\epsilon=\pm1$, and $F_p(\epsilon)\in\mathbb Z_p^\times$,
the exponent is $2s$ for every length $mp^s$ with $m,s\ge1$
(Part~II, Theorem~\ref{II:bdx:main}).
\resultheading{Supersingular Ap\'ery disks}{Part I, p.~\pageref{I:thm:aperydisks}}
For $p>3$, $p^2\mid N$, and the smoothness and supersingularity
conditions of Theorem~\ref{I:thm:aperydisks}, the two-step defect is
divisible by $p^{3v_p(N)-1}$.
\resultheading{The AESZ4 special value}{Parts III--IV, pp.~\pageref{paper:III}, \pageref{paper:IV}}
Part~III determines the cutoff residual and the generic residual monodromy.
Part~IV, Theorem~\ref{IV:thm:main}, gives $S_p\equiv a_p\pmod{p^3}$ for
$p\ge5$ at $t=-1/8$. The printed quantifier $p\ne3$ includes $p=2$,
where $S_2=161/162$ and $a_2=2$.
\par
\medskip
\noindent Section, equation and appendix numbers are local to each part.
All four parts share a bibliography. The new arguments in this research
manuscript await independent external review.
\clearpage
\tableofcontents
% Source: Beukers_Delaygue_and_Supersingular_Apery.tex
\begingroup
\newcommand{\Zp}{\mathbb Z_p}
\newcommand{\Fp}{F_p}
\newcommand{\eps}{\epsilon}
\newcommand{\al}{\alpha}
\newcommand{\ord}{\operatorname{ord}}
\openpaper{I}{Reflection Operators and Supersingular Ap\'ery Supercongruences}{Reflection operators and supersingular Ap\'ery truncations}{9 November 2026}{{\small Research manuscript prepared for Keio\par}\smallskip }




\begin{abstract}
For the truncations $F_L(\epsilon)=\sum_{n<L}((1/2)_n/n!)^d\epsilon^n$,
we prove the full mathematical statement of Beukers--Delaygue's
Conjecture~1.2 by a finite reflection calculation. Retaining this
reflection in an integral analytic digit operator gives the stronger,
all-length estimate
\[
 F_{mp^s}(\epsilon)-f_p(\epsilon)F_{mp^{s-1}}(\epsilon)
 \in p^{\min(2s,s+d-1)}\mathbb Z_p
\]
for every odd prime $p$, $d\ge2$, $\epsilon=\pm1$ with
$F_p(\epsilon)$ a unit, and $m,s\ge1$.
It proves the precision predicted by their Conjecture~1.5 throughout
$s\le d-1$, and Part~II, Theorem~\ref{II:bdx:main}, extends it to every higher level.
Separately, for the Ap\'ery polynomials $P_N(t)$, an integral modular
trace and an explicit elliptic crystal yield
\[
 P_N(t_0)-p^2P_{N/p^2}(t_0)
 \in p^{3v_p(N)-1}\mathbb Z_p.
\]
This holds for every $p^2\mid N$ on each smooth supersingular residue
disk specified in Theorem~\ref{I:thm:aperydisks}, for $p>3$.
The latter theorem includes the two rational CM values $t_0=\pm1$
under their respective inert-prime conditions and also non-CM lifts.
Complete proofs and exact finite verification programs are included.
\end{abstract}

\noindent\fbox{\begin{minipage}{0.94\textwidth}
\textbf{Scope and research status.}
Conjecture~1.2 of \cite{I:BD} is proved in its full stated range.
Conjecture~1.5 is proved here for $1\le s\le d-1$;
for $s\ge d$ the argument in this part gives $s+d-1$. The lift in Part~II gives $2s$ throughout. The Ap\'ery theorem has the explicit smoothness and
supersingularity hypotheses of Theorem~\ref{I:thm:aperydisks}.
The arguments have received internal independent checks; no completed
external peer review or claim of literature priority is asserted.
No unpublished BV3 or Swisher manuscript is assumed.
\end{minipage}}

\medskip
\noindent\textbf{Reading guide.}
Sections~1--4 prove the fixed-modulus reflection theorem.
Section~\ref{I:sec:transfer} supplies the stronger analytic operator and
its all-length precision bound. Section~\ref{I:sec:apery} gives the
independent Ap\'ery trace, integral jet comparison and supersingular
specialization. The appendices contain supplementary exact checks;
none is used as a substitute for a proof.

\clearpage

\section{Statements and relation to the published conjecture}

Throughout, $p$ is an odd prime and $d$ is an integer with $d\ge2$.
Set
\[
 \al_k=\frac{(1/2)_k}{k!}=\frac{\binom{2k}{k}}{4^k},\qquad
 F_L(z)=\sum_{k=0}^{L-1}\al_k^d z^k\quad(L\ge1).
\]
In particular, $F_1(z)=1$.  All coefficients belong to $\Zp$, since
their displayed denominators are powers of $4$.
Write
\begin{equation}\label{I:eq:sign}
 N=\frac{p-1}{2},\qquad
 \chi=\chi_{p,d}(\eps)=(-1)^{dN}\eps^N,
 \qquad \eps\in\{1,-1\}.
\end{equation}

\begin{theorem}[Block congruence]\label{I:thm:block}
If $\chi=1$, then for every integer $m\ge1$,
\begin{equation}\label{I:eq:block}
 F_{mp}(\eps)\equiv F_p(\eps)F_m(\eps)\pmod{p^2\Zp}.
\end{equation}
No unit assumption on $F_p(\eps)$ is required for this statement.
\end{theorem}

The sign condition holds for both values of $\eps$ if
$p\equiv1\pmod4$.  If $p\equiv3\pmod4$, it holds exactly when
$\eps=(-1)^d$.

\begin{theorem}[The Beukers--Delaygue unit-root congruence]\label{I:thm:unitroot}
Suppose $F_p(\eps)\in\Zp^\times$.  The Dwork limit
\begin{equation}\label{I:eq:limit}
 f_p(\eps)=\lim_{s\longrightarrow\infty}
       \frac{F_{p^s}(\eps)}{F_{p^{s-1}}(\eps)}
\end{equation}
exists, and
\begin{equation}\label{I:eq:unitroot}
 F_p(\eps)\equiv f_p(\eps)\pmod{p^2\Zp}.
\end{equation}
\end{theorem}

The existence of \eqref{I:eq:limit} is the published Dwork input recorded
in the introduction of \cite{I:BD}; it is not a new assertion of this
note.  The congruence \eqref{I:eq:unitroot} is exactly their
Conjecture~1.2.  Their Theorem~1.1 proves it at
$\eps_p=(-1)^{d(p-1)/2}$, and their Corollary~2.4 excludes the opposite
sign from the unit case when $p\equiv3\pmod4$.
They identify $p\equiv1\pmod4$, $\eps=-1$ as the remaining case.
The reflection calculation below applies to that case as well.

\section{Elementary congruences}

The calculations in this section also give self-contained proofs of the
elementary inputs appearing as Lemmas~2.1--2.5 of \cite{I:BD}.

\begin{lemma}[Babbage congruence]\label{I:lem:babbage}
For integers $A\ge B\ge0$,
\[
 \binom{pA}{pB}\equiv\binom AB\pmod{p^2}.
\]
\end{lemma}
\begin{proof}
The boundary cases are immediate.  In the product for
$\binom{pA}{pB}$, separate the factors indexed by multiples of $p$:
\[
 \binom{pA}{pB}
 =\binom AB
   \prod_{\substack{1\le j\le p(A-B)\\p\nmid j}}
                 \left(1+\frac{pB}{j}\right).
\]
The second factor is congruent modulo $p^2$ to
$1+pB\sum_{1\le j\le p(A-B),\,p\nmid j}j^{-1}$.
Each complete block contributes $\sum_{a=1}^{p-1}a^{-1}=0$ in
$\mathbb F_p$: inversion permutes the nonzero residues, whose sum is
zero for odd $p$.  The product is therefore $1$ modulo $p^2$.
\end{proof}

Put
\begin{equation}\label{I:eq:hgamma}
 h_t=\sum_{j=1}^{2t}\frac{(-1)^{j-1}}j\quad(0\le t\le N),
 \qquad \gamma=\frac{4^{p-1}-1}{p}\in\mathbb Z.
\end{equation}
The empty sum defining $h_0$ is zero.

\begin{lemma}[Alternating harmonic sum]\label{I:lem:eisenstein}
One has
\[
 \sum_{j=1}^{p-1}\frac{(-1)^{j-1}}j\equiv\gamma\pmod p.
\]
\end{lemma}
\begin{proof}
The identity
\[
 \frac{2^p-2}{p}
 =\sum_{j=1}^{p-1}\frac1j\binom{p-1}{j-1}
\]
and $\binom{p-1}{j-1}\equiv(-1)^{j-1}\pmod p$ show that the left
side of the desired congruence equals $(2^p-2)/p$ modulo $p$.
If $q=(2^{p-1}-1)/p$, then
$\gamma=q(2^{p-1}+1)\equiv2q=(2^p-2)/p\pmod p$.
\end{proof}

\begin{lemma}[Two reflections]\label{I:lem:reflections}
For $0\le t\le N$,
\begin{align}
 \al_{N-t}&\equiv(-1)^N\al_t\pmod p,\label{I:eq:alphareflect}\\
 h_{N-t}&\equiv\gamma-h_t\pmod p.\label{I:eq:hreflect}
\end{align}
Consequently, for $w_t=\al_t^d\eps^t$,
\begin{equation}\label{I:eq:wreflect}
 w_{N-t}\equiv\chi w_t\pmod p.
\end{equation}
\end{lemma}
\begin{proof}
Since $1/2\equiv-N\pmod p$ and $t!$ is a unit,
\[
 \al_t\equiv\frac{(-N)_t}{t!}
            =(-1)^t\binom Nt\pmod p.
\]
Binomial symmetry proves \eqref{I:eq:alphareflect}.
For \eqref{I:eq:hreflect}, subtract from the complete alternating sum
the final $2t$ terms.  Under $j=p-k$, those terms satisfy
\[
 \frac{(-1)^{j-1}}j
 =\frac{(-1)^{p-k-1}}{p-k}
 \equiv\frac{(-1)^{k-1}}k\pmod p.
\]
Their sum is therefore $h_t$ modulo $p$, and
Lemma~\ref{I:lem:eisenstein} supplies the complete sum.
Finally $\eps^{N-t}=\eps^N\eps^t$, giving \eqref{I:eq:wreflect}.
\end{proof}

\begin{lemma}[Uniform digit expansion]\label{I:lem:digits}
Let $a\ge0$ and $0\le t<p$ be integers.  If $t>N$, then
$\al_{pa+t}\in p\Zp$.  If $t\le N$, then
\begin{equation}\label{I:eq:digits}
 \al_{pa+t}\equiv
 \al_a\al_t\bigl(1+pa(2h_t-\gamma)\bigr)\pmod{p^2\Zp}.
\end{equation}
There is no upper bound on $a$ in this assertion.
\end{lemma}
\begin{proof}
The exact identity
\begin{equation}\label{I:eq:binomproduct}
 \binom{2pa+2t}{pa+t}
 =\binom{2pa}{pa}
   \frac{\prod_{j=1}^{2t}(2pa+j)}{\prod_{j=1}^t(pa+j)^2}
\end{equation}
has unit denominator in its second factor.  If $t>N$, its numerator
contains the factor $2pa+p$, proving the divisibility assertion.
If $t\le N$, all displayed $j$ are prime to $p$, and the second factor
is
\[
 \binom{2t}{t}
 \left(1+2pa\left(\sum_{j=1}^{2t}\frac1j
                         -\sum_{j=1}^{t}\frac1j\right)\right)
 \equiv\binom{2t}{t}(1+2pa h_t)\pmod{p^2}.
\]
Apply Lemma~\ref{I:lem:babbage} to the first binomial coefficient in
\eqref{I:eq:binomproduct}.  After division by $4^{pa+t}$, the remaining
normalization factor is
\[
 4^{a(1-p)}=(1+p\gamma)^{-a}\equiv1-pa\gamma\pmod{p^2}.
\]
Multiplying these expressions proves \eqref{I:eq:digits}.
\end{proof}

\section{Cancellation and the block theorem}

\begin{lemma}[Cancellation of the block correction]\label{I:lem:cancel}
If $\chi=1$, then
\begin{equation}\label{I:eq:cancel}
 \sum_{t=0}^N(2h_t-\gamma)\al_t^d\eps^t\equiv0\pmod p.
\end{equation}
If $F_p(\eps)$ is a $p$-adic unit, then necessarily $\chi=1$.
\end{lemma}
\begin{proof}
Let $S=\sum_{t=0}^Nh_tw_t$ and $A=\sum_{t=0}^Nw_t$.
When $\chi=1$, reflection and Lemma~\ref{I:lem:reflections} give
\[
 S=\sum_{t=0}^Nh_{N-t}w_{N-t}
   \equiv\sum_{t=0}^N(\gamma-h_t)w_t
   =\gamma A-S\pmod p.
\]
Thus $2S-\gamma A\equiv0$, which is \eqref{I:eq:cancel}.
For the second assertion, Lemma~\ref{I:lem:digits} with $a=0$ implies
\begin{equation}\label{I:eq:half}
 F_p(\eps)\equiv A\pmod{p^2},
\end{equation}
because $d\ge2$.  If $\chi=-1$, reflection gives $A\equiv-A\pmod p$.
Since $p$ is odd, this would imply $F_p(\eps)\equiv0\pmod p$.
\end{proof}

\begin{proof}[Proof of Theorem~\ref{I:thm:block}]
Write every summation index uniquely as $pa+t$, with $0\le a<m$ and
$0\le t<p$.  By Lemma~\ref{I:lem:digits}, the terms with $t>N$ vanish
modulo $p^2$ after taking the $d$th power.  The remaining terms satisfy
\[
 \al_{pa+t}^d\eps^{pa+t}
 \equiv\al_a^d\eps^a\al_t^d\eps^t
       \bigl(1+pda(2h_t-\gamma)\bigr)\pmod{p^2},
\]
where $\eps^{pa}=\eps^a$ because $p$ is odd.  Summing yields
\begin{align*}
 F_{mp}(\eps)
 &\equiv F_m(\eps)\sum_{t=0}^N\al_t^d\eps^t\\
 &\quad+pd\left(\sum_{a=0}^{m-1}a\al_a^d\eps^a\right)
          \left(\sum_{t=0}^N(2h_t-\gamma)\al_t^d\eps^t\right)
          \pmod{p^2}.
\end{align*}
The second line vanishes by Lemma~\ref{I:lem:cancel}, and the first
line is $F_m(\eps)F_p(\eps)$ by \eqref{I:eq:half}.
\end{proof}

\begin{corollary}\label{I:cor:levels}
If $\chi=1$, then for every $s\ge1$,
\[
 F_{p^s}(\eps)\equiv F_p(\eps)F_{p^{s-1}}(\eps)
              \equiv F_p(\eps)^s\pmod{p^2}.
\]
\end{corollary}
\begin{proof}
Use $m=p^{s-1}$ in Theorem~\ref{I:thm:block} and then induct on $s$.
\end{proof}

\section{Passage to the unit-root limit}

\begin{proof}[Proof of Theorem~\ref{I:thm:unitroot}]
By Lemma~\ref{I:lem:cancel}, the unit assumption implies $\chi=1$.
Corollary~\ref{I:cor:levels} shows that every $F_{p^s}(\eps)$ is a unit
and that, for $s\ge1$,
\begin{equation}\label{I:eq:quotients}
 R_s:=\frac{F_{p^s}(\eps)}{F_{p^{s-1}}(\eps)}
             \equiv F_p(\eps)\pmod{p^2}.
\end{equation}
The Dwork congruences as stated in \cite[Introduction, equation (1)]{I:BD}
give $R_{s+1}\equiv R_s\pmod{p^s}$; here $\eps^p=\eps$.
Hence the sequence $(R_s)$ converges in $\Zp$ to the value denoted by
$f_p(\eps)$ in \eqref{I:eq:limit}.
Since $p^2\Zp$ is closed, taking the limit in \eqref{I:eq:quotients}
proves \eqref{I:eq:unitroot}.
\end{proof}

\begin{remark}[What changes in the published proof]
The block expansion and the paired harmonic-sum calculation are the
method of \cite[Section 2]{I:BD}.  The additional observation is to retain
the joint factor $\chi=(-1)^{dN}\eps^N$ in the reflection, instead of
specializing in advance to $\eps=\eps_p$.
For $p\equiv1\pmod4$, the integer $N$ is even, so this factor is $1$
also at $\eps=-1$.  Thus the same cancellation completes the missing
sign without any new modular, crystalline, or unpublished input.
\end{remark}

\begin{remark}[Fixed precision versus growing precision]
The conclusion for every level in Corollary~\ref{I:cor:levels} has the
fixed modulus $p^2$.  It does not supply
$F_{p^s}(\eps)-f_p(\eps)F_{p^{s-1}}(\eps)\in p^{2s}\Zp$.
The latter assertion is Conjecture~1.5 of \cite{I:BD}.
Section~\ref{I:sec:transfer} proves it for $s\le d-1$ and obtains
the weaker exponent $s+d-1$ for $s\ge d$.
We do not identify $f_p(\eps)$ with a particular modular-form
coefficient or an explicitly specified motivic Euler factor.
\end{remark}

For a concrete instance of the added sign, take $p=5$, $d=2$, and
$\eps=-1$.  Direct rational arithmetic gives
\[
 F_5(-1)=1-\frac14+\frac9{64}-\frac{25}{256}
                 +\frac{1225}{16384}
          =\frac{14217}{16384}\equiv13\pmod{25}.
\]
This is a unit, so $f_5(-1)\equiv13\pmod{25}$ by
Theorem~\ref{I:thm:unitroot}.  The block theorem also gives
$F_{25}(-1)\equiv13^2\equiv19\pmod{25}$.


\section{A precision bound for every level and every length}
\label{I:sec:transfer}

The reflection can also be retained in an integral analytic operator.
Keeping track of its even and odd parts yields a stronger result than
the fixed-modulus congruence.
Put
\[
 \nu_d(s)=\min(2s,s+d-1)=s+\min(s,d-1).
\]

\begin{theorem}[All-length precision bound]\label{I:thm:alllength}
Let $p$ be any odd prime, $d\ge2$, and $\eps\in\{1,-1\}$.
Suppose $F_p(\eps)\in\Zp^\times$, and let $f_p(\eps)$ be the limit
in \eqref{I:eq:limit}.  For all integers $m,s\ge1$,
\begin{equation}\label{I:eq:alllength}
 F_{mp^s}(\eps)-f_p(\eps)F_{mp^{s-1}}(\eps)
                  \in p^{\nu_d(s)}\Zp.
\end{equation}
No unit assumption on $F_m(\eps)$ is needed.
Equivalently, if $K\ge1$ and $k=v_p(K)\ge1$, then
\begin{equation}\label{I:eq:cutoffvaluation}
 F_K(\eps)-f_p(\eps)F_{K/p}(\eps)
                  \in p^{k+\min(k,d-1)}\Zp.
\end{equation}
\end{theorem}

\begin{corollary}[A range of the growing-precision conjecture]
\label{I:cor:bd15range}
Under these assumptions, for every $1\le s\le d-1$,
\[
 F_{p^s}(\eps)\equiv f_p(\eps)F_{p^{s-1}}(\eps)
                                     \pmod{p^{2s}}.
\]
Thus the $s$th assertion of Beukers--Delaygue's Conjecture~1.5 is
proved for all $d\ge s+1$.  At levels $s\ge d$,
Theorem~\ref{I:thm:alllength} instead supplies the modulus $p^{s+d-1}$.
\end{corollary}

The proof below is elementary.  In particular, the exact reflection of
the interpolating functions is proved from finite products; no property
of the $p$-adic gamma function is required.

\subsection{Integral analytic interpolation of digit blocks}

Let $\mathcal A=\Zp\langle a\rangle$ be the algebra of restricted power
series $\sum_{j\ge0}b_ja^j$, where $b_j\in\Zp$ and $b_j\to0$.
We use its Gauss norm.  Divisibility in $\mathcal A$ always means
divisibility of every power-series coefficient.  Integral affine
substitution, evaluation at a point of $\Zp$, and finite integral sums
are bounded operations on this algebra.

\begin{lemma}[Analytic digit functions]\label{I:lem:analyticdigits}
There are functions $g_t\in\mathcal A$, $0\le t<p$, such that for
every nonnegative integer $a$,
\begin{equation}\label{I:eq:exactdigits}
 g_t(a)=\frac{\al_{pa+t}}{\al_a}.
\end{equation}
For $0\le t\le N$, there is $b_t\in\Zp$ with
\begin{equation}\label{I:eq:affinegt}
 g_t(a)\equiv\al_t(1+pb_ta)\pmod{p^2\mathcal A}.
\end{equation}
For $N<t<p$, one has $g_t\in p\mathcal A$.
All these statements hold also for $p=3$.
\end{lemma}
\begin{proof}
First interpolate $g_0(a)=\al_{pa}/\al_a$.
For nonnegative integer $a$, cancellation of the numerator and
denominator factors divisible by $p$ gives
\[
 \frac{g_0(a+1)}{g_0(a)}=R(a),\qquad
 R(a)=C\,
 \frac{\displaystyle\prod_{\substack{0\le t<p\\t\ne N}}
                 (1+pa/(t+1/2))}
      {\displaystyle\prod_{i=1}^{p-1}(1+pa/i)},
 \qquad
 C=\frac{\displaystyle\prod_{\substack{0\le t<p\\t\ne N}}(t+1/2)}
          {\displaystyle\prod_{i=1}^{p-1}i}.
\]
Both constant-factor lists reduce to $\mathbb F_p^\times$, so
$C\in1+p\Zp$.  For $j\ge1$ put
\[
 s_j=\sum_{\substack{0\le t<p\\t\ne N}}(t+1/2)^{-j}
                  -\sum_{i=1}^{p-1}i^{-j}\in p\Zp.
\]
The logarithm, defined by its convergent power series on $1+p\mathcal A$,
is
\begin{equation}\label{I:eq:logR}
 \log R(a)=\log C+
       \sum_{j\ge1}\frac{(-1)^{j-1}p^js_j}{j}a^j.
\end{equation}
The polynomial $P_j(a)=\sum_{u=0}^{a-1}u^j$, initially defined on
nonnegative integers, has the expression
\[
 P_j(a)=\sum_{v=0}^j
 \left\{\begin{matrix}j\\v\end{matrix}\right\}
 \frac{a(a-1)\cdots(a-v)}{v+1}.
\]
Here the braces denote Stirling numbers of the second kind.  Consequently
every power-basis coefficient of $P_j$ has valuation at least
$-\lfloor\log_p(j+1)\rfloor$.
Define
\begin{equation}\label{I:eq:logg}
 L(a)=a\log C+
       \sum_{j\ge1}\frac{(-1)^{j-1}p^js_j}{j}P_j(a).
\end{equation}
The coefficients contributed by its $j$th summand have valuation at
least
\begin{equation}\label{I:eq:coefficientbound}
 j+1-v_p(j)-\lfloor\log_p(j+1)\rfloor\ge2,
\end{equation}
and this lower bound tends to infinity.  For completeness, the inequality
is immediate for $j=1,2,3$.  For $j\ge4$ use
$p\ge3$ and $3^{j-1}\ge j(j+1)$; the latter holds at $j=4$ and follows
thereafter by induction.
Thus \eqref{I:eq:logg} converges in $\mathcal A$,
$L\in p\mathcal A$, and $L-a\log C\in p^2\mathcal A$.

Since $p$ is odd, the exponential converges on $p\mathcal A$.  Set
$g_0=\exp L$.  The identities $P_j(a+1)-P_j(a)=a^j$ give
$g_0(a+1)=R(a)g_0(a)$ and $g_0(0)=1$.  These identify its nonnegative
integer values with $\al_{pa}/\al_a$.  Moreover,
$g_0(a)\equiv1+a\log C\pmod{p^2\mathcal A}$.
For the remaining digits set
\begin{equation}\label{I:eq:gtproduct}
 g_t(a)=g_0(a)4^{-t}
       \frac{\prod_{j=1}^{2t}(2pa+j)}
            {\prod_{j=1}^t(pa+j)^2}.
\end{equation}
All denominator constants are units, so this expression lies in
$\mathcal A$.  The exact binomial product \eqref{I:eq:binomproduct}
proves \eqref{I:eq:exactdigits}.  If $t\le N$, expansion of the unit
factors proves \eqref{I:eq:affinegt}, with constant term $\al_t$.
If $t>N$, the numerator has the factor $p(2a+1)$, proving
$g_t\in p\mathcal A$.
\end{proof}

\subsection{An exact reflection from finite products}

\begin{lemma}[Exact analytic reflection]\label{I:lem:exactreflection}
Let $r(a)=-a-1/2$.  For every $0\le t\le N$,
\begin{equation}\label{I:eq:exactreflection}
 g_{N-t}(r(a))=(-1)^N g_t(a)
                                      \qquad\text{in }\mathcal A.
\end{equation}
\end{lemma}
\begin{proof}
Put $N_j=(p^j-1)/2$.  For each fixed nonnegative integer $n$ and all
sufficiently large $j$, a finite product gives
\[
 \frac{\al_{N_j-n}}{\al_{N_j}}
       =\prod_{k=0}^{n-1}\frac{N_j-k}{N_j-k-1/2}
       \longrightarrow\prod_{k=0}^{n-1}\frac{k+1/2}{k+1}
       =\al_n.
\]
The convergence follows from $N_j\to-1/2$ in $\Zp$; every limiting
denominator is nonzero in $\mathbb Q_p$.
Also, since $N_j=pN_{j-1}+N$,
\[
 \frac{\al_{N_j}}{\al_{N_{j-1}}}
     =g_N(N_{j-1})\longrightarrow C_*:=g_N(-1/2).
\]
Fix a nonnegative integer $a$ and $0\le t\le N$, and set
$n=pa+t$, $a_j=N_{j-1}-a$.  For large $j$,
$pa_j+(N-t)=N_j-n$.  Applying the two finite-product limits above
to the numerator and denominator of
$g_{N-t}(a_j)=\al_{N_j-n}/\al_{N_{j-1}-a}$ yields
\[
 g_{N-t}(r(a))=C_*\frac{\al_n}{\al_a}=C_*g_t(a).
\]
Continuity extends this equality from nonnegative integers to $\Zp$.
It is also an identity of restricted power series: a nonzero convergent
power series cannot vanish at all $p^j$, which approach zero.
Applying the identity twice gives $C_*^2=1$, using $g_0(0)=1$.
By \eqref{I:eq:affinegt} and \eqref{I:eq:alphareflect},
$C_*\equiv\al_N\equiv(-1)^N\pmod p$.
The two roots $1,-1$ are distinct modulo odd $p$, so
$C_*=(-1)^N$.
\end{proof}

\subsection{Digit operators and their two reflection components}

Define bounded integral operators on $\mathcal A$ by
\begin{align}
 (Sf)(a)&=\sum_{t=0}^{p-1}\eps^t g_t(a)^d f(pa+t),
                                                      \label{I:eq:fulloperator}\\
 (T_0f)(a)&=\sum_{t=0}^{N}\eps^t g_t(a)^d f(pa+t).
                                                      \label{I:eq:loweroperator}
\end{align}
By Lemma~\ref{I:lem:analyticdigits},
\begin{equation}\label{I:eq:perturbation}
 S=T_0+p^dK
\end{equation}
for a bounded integral operator $K$.
Let $\mathcal Rf=f\circ r$.  Throughout the remainder of the proof,
the ordinary hypothesis and Lemma~\ref{I:lem:cancel} give $\chi=1$.
Equation \eqref{I:eq:exactreflection}, together with
$r(pa+t)=pr(a)+(N-t)$, then gives
\begin{equation}\label{I:eq:commutation}
 T_0\mathcal R=\mathcal RT_0.
\end{equation}

\begin{lemma}[Operator divisibilities]\label{I:lem:operatorblocks}
Set $c=-1/4$, the fixed point of $r$, and $\ell(f)=f(c)$.
Decompose
\[
 \mathcal A=\Zp\cdot1\ \oplus\ E\ \oplus\ O,
\]
where $E=\{f:\mathcal Rf=f,\ \ell(f)=0\}$ and
$O=\{f:\mathcal Rf=-f\}$.
In this decomposition the block valuations of $S$ are at least
\begin{equation}\label{I:eq:blockmatrix}
 \begin{pmatrix}
 0&0&d\\
 2&2&d\\
 d&d&1
 \end{pmatrix}.
\end{equation}
An entry $b$ means that the corresponding component map has image
divisible by $p^b$.  The scalar entry $\ell(S1)$ is a unit.
\end{lemma}
\begin{proof}
The projections $(1\pm\mathcal R)/2$ are bounded and integral.
Every odd function vanishes at $c$, so this is indeed a bounded direct
sum decomposition.
For arbitrary $f\in\mathcal A$, the congruence
\[
 f(pa+t)\equiv f(t)+pa f'(t)\pmod{p^2\mathcal A}
\]
and \eqref{I:eq:affinegt} show that $T_0f$ is affine modulo $p^2$.
If $f$ is even, \eqref{I:eq:commutation} makes $T_0f$ even.  An affine
polynomial $A+Ba$ fixed by $a\mapsto-a-1/2$ modulo $p^2$ has
$B\in p^2\Zp$.  Thus its nonconstant even component is divisible by
$p^2$.  If $f$ is odd, $T_0f$ is odd and is constant modulo $p$;
that constant must be zero, so $T_0f\in pO$.
These observations give \eqref{I:eq:blockmatrix} for $T_0$, with zero
off-diagonal reflection components.  Equation \eqref{I:eq:perturbation}
gives the displayed matrix for $S$, since $d\ge2$.
Finally $S1$ is constant modulo $p$ with value
$\sum_{t=0}^N\al_t^d\eps^t\equiv F_p(\eps)$, proving the unit assertion.
\end{proof}

\subsection{Normalized iteration and its precision}

\begin{lemma}[Iteration estimate]\label{I:lem:iteration}
There is a unit $\lambda\in\Zp^\times$ such that, for all $s\ge1$,
\begin{equation}\label{I:eq:functioncongruence}
 S^s1-\lambda S^{s-1}1\in p^{\nu_d(s)}\mathcal A.
\end{equation}
\end{lemma}
\begin{proof}
On the set $\mathcal Q=1+p^2E+p^dO$ define
\[
 \mathcal N(h)=\frac{Sh}{\ell(Sh)}.
\]
The denominator is a unit by \eqref{I:eq:blockmatrix}, and the same
matrix shows that $\mathcal N$ maps $\mathcal Q$ into itself.
Let $h_0=1$, $h_s=\mathcal N(h_{s-1})$, and
$Z_s=\ell(S^s1)$.  Induction gives $Z_s\in\Zp^\times$ and
$h_s=S^s1/Z_s$.

For $g,h\in\mathcal Q$ the exact identity
\begin{equation}\label{I:eq:normalizeddifference}
 \mathcal N(g)-\mathcal N(h)
 =\frac{S(g-h)-\ell(S(g-h))\mathcal N(h)}{\ell(Sg)}
\end{equation}
shows the following.  If the even and odd parts of $g-h$ belong to
$p^uE$ and $p^vO$, respectively, then the parts of the left side belong
to
\[
 p^{\min(u+2,v+d)}E,
 \qquad p^{\min(u+d,v+1)}O.
\]
Indeed, \eqref{I:eq:blockmatrix} gives these bounds for $S(g-h)$;
the subtracted scalar has valuation at least $\min(u,v+d)$ and
multiplies even and odd components divisible by $p^2$ and $p^d$.

Put $\Delta_s=h_s-h_{s-1}$.  Starting with
$\Delta_1\in p^2E+p^dO$, induction in
\eqref{I:eq:normalizeddifference} gives
\begin{equation}\label{I:eq:incrementbounds}
 \Delta_s\in p^{u_s}E+p^{v_s}O,
 \qquad
 u_s=\min(2s,2d+s-2),\quad v_s=d+s-1.
\end{equation}
For example, the inductive lower bounds are
$\min(u_{s-1}+2,v_{s-1}+d)=u_s$ and
$\min(u_{s-1}+d,v_{s-1}+1)=v_s$ for $s\ge2$.
Thus $h_s$ converges in $\mathcal A$ to $h_\infty\in\mathcal Q$.

Set
\[
 z_s=\frac{Z_s}{Z_{s-1}}=\ell(Sh_{s-1}),
 \qquad \lambda=\ell(Sh_\infty)\in\Zp^\times.
\]
The scalar row of \eqref{I:eq:blockmatrix} gains an additional $d$
powers of $p$ on odd inputs.  Since
$h_\infty-h_{s-1}=\sum_{j=s}^{\infty}\Delta_j$, it follows that
\begin{equation}\label{I:eq:centeredratio}
 \lambda-z_s\in p^{u_s}\Zp.
\end{equation}
Here $u_j$ is nondecreasing and
$\min(u_j,v_j+d)=u_j$.
Consequently,
\[
 S^s1-\lambda S^{s-1}1
 =Z_{s-1}\bigl(z_s(h_s-h_{s-1})+(z_s-\lambda)h_{s-1}\bigr).
\]
Equations \eqref{I:eq:incrementbounds}--\eqref{I:eq:centeredratio}
give valuation at least
$\min(u_s,v_s)=\min(2s,s+d-1)=\nu_d(s)$ throughout this function.
This proves \eqref{I:eq:functioncongruence}.
\end{proof}

\subsection{Evaluation at the original truncations}

\begin{proof}[Proof of Theorem~\ref{I:thm:alllength}]
For $m\ge1$ define the bounded integral functional
\[
 \mathcal L_m(f)=\sum_{a=0}^{m-1}\al_a^d\eps^a f(a).
\]
Exact digit expansion \eqref{I:eq:exactdigits} gives
\begin{equation}\label{I:eq:prefixidentity}
 \mathcal L_m(S^s1)=F_{mp^s}(\eps)\qquad(s\ge0).
\end{equation}
To verify this, an iterated digit sequence represents a unique
$b$ with $0\le b<p^s$.  The product of its $g_t$ factors telescopes
to $\al_{p^sa+b}/\al_a$.  Its sign, after multiplication by
$\eps^a$, is $\eps^{p^sa+b}$, because $p$ is odd.
Summing over $a$ and $b$ gives \eqref{I:eq:prefixidentity}.
This identity uses rational equalities and does not require
$\al_a$ or $F_m(\eps)$ to be a unit.

For $m=1$, $F_{p^s}(\eps)=(S^s1)(0)=Z_sh_s(0)$.
Every $h_s(0)$ is a unit since $h_s\in1+p^2E+p^dO$.
Therefore \eqref{I:eq:functioncongruence}, evaluated at zero, implies
\[
 \frac{F_{p^s}(\eps)}{F_{p^{s-1}}(\eps)}
       \equiv\lambda\pmod{p^{\nu_d(s)}}.
\]
Since $\nu_d(s)\to\infty$, the limit equals $\lambda$.
It is precisely the limit $f_p(\eps)$ defined in
\eqref{I:eq:limit}.  In particular, this operator argument also proves
the convergence needed here.
Applying $\mathcal L_m$ to \eqref{I:eq:functioncongruence} now proves
\eqref{I:eq:alllength}, with the same $\lambda=f_p(\eps)$ for every $m$.
Finally, write $K=mp^k$ with $p\nmid m$ to obtain
\eqref{I:eq:cutoffvaluation}.
\end{proof}

\begin{remark}[The original evaluation point]
Normalization took place at $c=-1/4$ only to separate reflection-even
and reflection-odd functions.  The sums themselves are evaluated at
$a=0$, as \eqref{I:eq:prefixidentity} shows.  In particular,
\[
 \frac{F_{p^s}(\eps)}{F_{p^{s-1}}(\eps)}
     =z_s\frac{h_s(0)}{h_{s-1}(0)}.
\]
The second quotient differs from $1$ by an element of
$p^{\nu_d(s)}\Zp$ by \eqref{I:eq:incrementbounds}, while
$z_s-\lambda\in p^{u_s}\Zp$ by \eqref{I:eq:centeredratio}.
Thus centering the normalization does not assert an unjustified
identification of evaluations at $c$ and at zero.
\end{remark}

\begin{remark}[Scope of the higher precision]
Under the unit hypothesis, the reflection estimate gives the precision of
Conjecture~1.2 and the exponent $2s$ of Conjecture~1.5 for $1\le s\le d-1$.
For $s\ge d$ it gives $s+d-1$. The lifted reflection of Part~II,
Theorem~\ref{II:bdx:main}, gives $2s$ for these levels as well.
\end{remark}


\section{Ap\'ery truncations on supersingular residue disks}\label{I:sec:apery}

In this section the notation $F$ refers to the Ap\'ery period. Put
\[
 A_n=\sum_{k=0}^n\binom nk^2\binom{n+k}k^2,\qquad
 F(t)=\sum_{n\geq0}A_nt^n,\qquad P_K(t)=\sum_{0\leq n<K}A_nt^n.
\]
We prove a congruence for every integral point in each smooth
supersingular residue disk of the Ap\'ery family. The proof first
establishes an all-cutoff formal congruence, then identifies an explicit
integral rank-two elliptic crystal, and finally specializes a transported
rank-three jet identity. Its crystalline inputs are standard;
the trace induction, differential residual, integral gauge and precision
comparison are proved here. The statement is an all-level extension of
first-level modular supercongruences, not a claim that the older Ap\'ery
coefficient congruences were previously open.

\subsection{The exact full-cutoff residual}

Set
\[
 L=\theta^3-t(2\theta+1)(17\theta^2+17\theta+5)
       +t^2(\theta+1)^3,\qquad\theta=t\frac d{dt}.
\]
The Ap\'ery recurrence gives the exact identity
\begin{equation}\label{I:ap:residual}
 LP_K=K^3t^K(-A_K+tA_{K-1}).
\end{equation}
Indeed the coefficients below $K$ vanish. At degree $K$ the recurrence
would have supplied $K^3 A_K$, and at degree $K+1$ the remaining term is
$K^3 A_{K-1}$. There are no later terms. This proves
\eqref{I:ap:residual} over $\mathbb Z[t]$, not merely modulo a prime power.

\subsection{An independent integral trace proof}

The published modular parametrization in Beukers, arXiv:2403.03301v3,
Appendix C, Section 9.5 \cite{III:Beukers2026}, is
\[
 t=\left(\frac{\eta_1\eta_6}{\eta_2\eta_3}\right)^{12},\qquad
 \Theta=\frac{(\eta_2\eta_3)^7}{(\eta_1\eta_6)^5},\qquad F(t(q))=\Theta(q).
\]
Here $\eta_j=\eta(j\tau)$, $q=e^{2\pi i\tau}$, $t\in q+q^2\mathbb Z[[q]]$, $\Theta\in1+q\mathbb Z[[q]]$,
and the group is $G=\Gamma_0(6)+\langle W_6\rangle$. On $X_0(6)$,
the cusp orders, in denominator order $1,2,3,6$, are
\[
\begin{array}{c|rrrr}
c&1&2&3&6\\\hline
\operatorname{ord}_c t&1&-1&-1&1\\
\operatorname{ord}_c\Theta&0&1&1&0\\
\text{width}&6&3&2&1
\end{array}
\]
The eta transformation formula gives $t\circ W_6=t$ and
$\Theta|_2 W_6=-\Theta$. The two pole cusps are exchanged by $W_6$.
Since $t:X_0(6)\to\mathbb P^1$ has degree two, it induces a degree-one
function on $X(G)$.

For $p>3$, the integral determinant-$p$ matrices with lower left entry
divisible by six form the usual Hecke set. Every such matrix preserves
the denominator class of a cusp modulo six. To see this prime by prime,
reduce the matrix modulo $2$ or $3$: it is invertible and upper triangular,
so it preserves whether the denominator is divisible by that prime.
Cancelling a common factor can only remove a factor of $p$, which changes
neither condition. Thus each of the two pole cusp classes is preserved.
The Fricke normalizer preserves the Hecke set. Consequently the usual
trace construction applies, with the same cusp width at the source and
target of each branch.

For a polynomial $P\in\mathbb Z[T]$, put
\[
 Q_P(q)=\frac{\Theta(q)P(t(q^p))}{\Theta(q^p)},\qquad
 G_P(q)=\frac{\Theta(q^p)P(t(q))}{\Theta(q)}.
\]
With determinant-normalized weight-two slash operators, the sum of
$p\Theta(\tau)P(t(\beta\tau))/(\Theta|_2\beta)(\tau)$ over the
$p+1$ Hecke representatives equals
\begin{equation}\label{I:ap:trace}
 \operatorname{Tr}_p(P)=Q_P+p^3 U_pG_P.
\end{equation}
The finite $p$ representatives each contribute a factor $p^2$, and
root-of-unity summation contributes the remaining $p$.

The trace is $G$-invariant. It has no poles away from the pole orbit of
$t$, since $\Theta$ has no zeros there and the Hecke correspondence
preserves that orbit in both directions. At a pole cusp, one branch
multiplies local orders by $p$ and the other $p$ branches divide them by
$p$. The respective orders of $\Theta,t$ are $1,-1$. It follows that
\begin{equation}\label{I:ap:tracepoly}
 \operatorname{Tr}_p(P)=\Pi_P(t),\qquad
 \Pi_P\in\mathbb Z[T],\qquad \deg\Pi_P\leq p\deg P+p-1.
\end{equation}
The coefficients are integral because \eqref{I:ap:trace} has an integral
$q$-expansion and $t=q+O(q^2)$ has an integral inverse series.

Let $\sigma_p(t)=t(q(t)^p)$ and $\rho_p=F(t)/F(\sigma_p(t))$.
Then, for all $m,s\geq1$,
\begin{equation}\label{I:ap:formal}
 \rho_p(t)P_{mp^{s-1}}(\sigma_p(t))-P_{mp^s}(t)
 \in p^{3s}\mathbb Z_p[[t]].
\end{equation}
Here is the complete induction. Apply \eqref{I:ap:trace} to
$P=P_{mp^{s-1}}$. Its trace polynomial has degree $<mp^s$.
By the preceding-level assertion,
\[
 G_P\equiv P_{mp^{s-2}}(t(q^p))\pmod{p^{3(s-1)}}.
\]
Thus $p^3U_pG_P\equiv p^3P_{mp^{s-2}}(t(q))\pmod{p^{3s}}$.
For $s=1$ interpret the preceding polynomial as zero; the statement
modulo $p^3$ is then automatic. Hence $Q_P$ is congruent modulo $p^{3s}$
to a polynomial of degree $<mp^s$. But $Q_P-F(t)$ has order at least
$mp^s$ in $t$. Comparing the initial $mp^s$ coefficients identifies that
polynomial with $P_{mp^s}$ and proves \eqref{I:ap:formal}.

\subsection{An explicit smooth elliptic cover and integral gauge}

Set
\[
 D=(1+8h)(1+9h),\qquad W=1-72h^2,\qquad t=\frac hD.
\]
Consider the elliptic curve
\begin{equation}\label{I:ap:elliptic}
 E_h:\quad y^2=x^3+\frac{1+12h+24h^2}{4}x^2+h^3(1+9h)x.
\end{equation}
Its discriminant and $c_4$ are
\[
 \Delta=h^6(1+8h)^3(1+9h)^2,\qquad
 c_4=(1+6h)(1+18h+84h^2+24h^3).
\]
Thus it is smooth over $\mathbb Z_p[h,1/(hD)]$ for every $p>3$.
With $\omega=dx/(2y)$ and $\eta=x\omega$, elementary differential
reduction gives
\begin{equation}\label{I:ap:GM}
 \nabla_h\binom\omega\eta=
 \begin{pmatrix}
 -6/(1+8h)&6/(hD)\\
 -6h^2/(1+8h)&6/(1+8h)
 \end{pmatrix}\binom\omega\eta.
\end{equation}
For example, substitute $f=x^3+ax^2+bx$ and reduce
$-f_h dx/(4y^3)$ modulo exact differentials $d(R/y)$. This is an
identity of relative de Rham cohomology, with denominators only those
displayed. In particular $(\omega,\nabla_h\omega)$ is an integral basis:
its determinant relative to $(\omega,\eta)$ is $6/(hD)$, a unit.

The normalized vanishing-cycle period $\ell$ of $\omega$ at $h=0$
has constant term one; the tangent lines at the node are $y=\pm x/2$.
Equation \eqref{I:ap:GM} gives
\[
 \ell''+\frac{1+34h+216h^2}{hD}\ell'
      +\frac{6(1+12h)}{hD}\ell=0.
\]
Hence $\ell=1-6h+42h^2-312h^3+\cdots$, the twisted Zagier F period.
The published eta parametrization is
\[
 h(q)=\frac{\eta_6^5\eta_2}{\eta_1^5\eta_3},\qquad
 \ell(h(q))=\frac{\eta_1^6\eta_6}{\eta_2^3\eta_3^2}.
\]
Logarithmic eta differentiation gives the useful exact identity
\[
 q\frac d{dq}\log h
 =\frac{-5E_2(\tau)+2E_2(2\tau)-3E_2(3\tau)+30E_2(6\tau)}{24}
 =\Theta=D\ell^2.
\]
The last equality can also be checked in the weight-two modular form
space for $\Gamma_0(6)$ by its coefficients through $q^2$. The
elliptic Wronskian and the multiplicative fiber of order six give
$d\log q_E/dh=6/(hD\ell^2)$, while $\Delta\sim h^6$, $c_4\sim1$
fix the leading constant. Consequently the Tate parameter of
\eqref{I:ap:elliptic} is exactly $q_E=q^6$. Canonical quotient of the
elliptic curve sends $q_E$ to $q_E^p$ and therefore sends $h(q)$ to
$h(q^p)$. This identifies the elliptic canonical lift with the modular
lift used in the scalar trace argument.
The eta identity, or a direct comparison of differential equations and
initial terms, gives
\begin{equation}\label{I:ap:square}
 F(t(h))=D(h)\ell(h)^2.
\end{equation}

No square root needs to be adjoined to construct the integral crystal.
Use instead the quadratic twist
\[
 E_h^D:\quad Y^2=D(h)
 \left(X^3+\frac{1+12h+24h^2}{4}X^2+h^3(1+9h)X\right),
 \qquad \omega_*=\frac{D(h)dX}{2Y}.
\]
Its period is $a(h)=\sqrt D\,\ell(h)$, with the formal square root
equal to one at $h=0$, and $a^2=F(t(h))$. As $D$ is a unit,
$(\omega_*,\nabla_h\omega_*)$ remains an integral basis. Explicitly,
the period equation simplifies to
\begin{equation}\label{I:ap:ranktwo}
 a''+\mathcal P a'+\mathcal Q a=0,\qquad
 \mathcal P=\frac1h,\quad
 \mathcal Q=-\frac{720h^2+169h+10}{4hD^2}.
\end{equation}

Fix $m\geq1$ with $(m,p)=1$, put $K_s=mp^s$, and set
$f_s(h)=P_{K_s}(t(h))$. Let $\mathcal L$ be the monic symmetric
square of \eqref{I:ap:ranktwo}. Since
\[
 1-34t+t^2=\frac{W^2}{D^2},\qquad t'=\frac W{D^2},
\]
the exact residual \eqref{I:ap:residual} becomes
\begin{equation}\label{I:ap:pullres}
 \mathcal L f_s
  =K_s^3\frac{W}{h^3D}\,t^{K_s}
               (-A_{K_s}+tA_{K_s-1}).
\end{equation}
This is integral on the whole smooth $h$-chart. In particular, no
inverse of $W$ is needed, despite the ramification of $h\mapsto t$.

\subsection{The integral Frobenius--Taylor comparison}

Write $A=\widehat{\mathbb Z_p[h,1/(hD)]}$. Let $\phi$ be any integral
Frobenius lift of this ring. In the integral basis
$(\omega_*,\nabla_h\omega_*)$, use the row convention
\[
 F_\phi(\phi^*\omega_i)=\sum_j(M_\phi)_{ij}\omega_j,
 \qquad V_\phi=pM_\phi^{-1}.
\]
Crystalline Frobenius and Verschiebung for a smooth elliptic family imply
that both $M_\phi,V_\phi$ are integral; their determinant is $p$ times
a unit. More explicitly, the cup product of
$(\omega_*,\nabla_h\omega_*)$ is $6/h$ up to the fixed sign convention,
so compatibility of Frobenius with cup product gives
\[
 \det M_\phi=p\frac{h}{\phi(h)}.
\]
Thus the integrality of $V_\phi$ also follows directly from the adjugate
of $M_\phi$. The integral comparison and Frobenius are those for the
Dieudonn\'e crystal of the elliptic scheme's $p$-divisible group:
Berthelot--Messing, \emph{Th\'eorie de Dieudonn\'e cristalline I} \cite{I:BM},
Section 2, Theorem 2 and the constructions on printed pp.\ 24--25;
Section 3, equation (3.1), identifies this crystal with relative
degree-one cohomology and its Hodge filtration.
The square-period jet matrix is
\[
 S=\begin{pmatrix}1&0&0\\0&2&0\\-2\mathcal Q&-2\mathcal P&2\end{pmatrix},
 \qquad B_\phi=S\operatorname{Sym}^2(V_\phi)(S^\phi)^{-1}.
\]
Its determinant-changing matrices are integral units, since $p>3$.
For $Y_s=(f_s,f_s',f_s'')^{\mathsf T}$ one has
\begin{equation}\label{I:ap:jet}
 Y_s-B_\phi Y_{s-1}^\phi\in p^{3s}A^3\qquad(s\geq1).
\end{equation}

Here is the proof, including the change-of-lift precision.
On the ordinary localization $A_H$, the canonical lift $\sigma$ has
$h(q)\mapsto h(q^p)$. The split Tate degeneration of
\eqref{I:ap:elliptic}, and $D(0)=1$, normalize the Hodge coefficient as
$c=p/\alpha$ with $\alpha=a(h)/a(h^\sigma)$. Horizontality then gives
\[
 M_\sigma=\begin{pmatrix}c&0\\c'/\sigma'&c/\sigma'\end{pmatrix},
 \quad V_\sigma=\begin{pmatrix}\alpha&0\\\alpha'&\alpha\sigma'\end{pmatrix}.
\]
Therefore $B_\sigma$ is exactly the chain-rule matrix obtained by
differentiating $a(h)^2/a(h^\sigma)^2$ twice. Its three columns are
divisible by $1,p,p^2$, respectively. The scalar formal congruence
\eqref{I:ap:formal}, together with its first two derivatives, gives
\eqref{I:ap:jet} for $\sigma$ on the ordinary chart. To pass from the
formal cusp identity to $A_H$, use the $q$-expansion principle modulo
$p^e$: the mod-$p$ special fiber is integral, hence localization and
expansion into the Laurent series field at $h=0$ are injective, and
successive division by $p$ proves saturation at every precision.
For the underlying ordinary canonical quotient and the normalized
Tate differential, see Katz, \emph{$p$-adic properties of modular
schemes and modular forms}, Sections 3.1--3.4 and Appendix A2,
especially (A2.2.1) and (A2.2.7) \cite{I:Katz}. Its Proposition 2.7.1 gives the
integral $q$-expansion principle. The equality $q_E=q^6$ above checks
that these standard constructions have exactly the normalization
used here.

Now $\delta=\phi(h)-\sigma(h)$ is divisible by $p$. Let
$T(\sigma,\phi)$ send a horizontal jet at $\phi(h)$ to its horizontal
continuation at $\sigma(h)$. The crystalline change-of-lift identity,
with the row convention for $M$, is
\[
 B_\phi=B_\sigma T(\sigma,\phi).
\]
This is the Taylor comparison of the same crystal at two lifts; see
the explicit divided-power formula in Stacks Project,
Proposition 60.17.4, Tag 07JH \cite{I:Stacks}.
Horizontal Taylor transport is integral: its terms have the form
$\delta^n/n!$ times iterated integral connection matrices, and
$n-v_p(n!)\to\infty$. The residual \eqref{I:ap:pullres} shows that
\[
 T(\sigma,\phi)Y_{s-1}^\phi-Y_{s-1}^\sigma
 \in p^{3(s-1)}\left(p^3A_H\oplus p^2A_H\oplus pA_H\right).
\]
Indeed, solve the homogeneous initial-value problem having the same
three-jet at the starting point. The difference from $f_{s-1}$ has
zero initial jet and forcing term in $p^{3(s-1)}A_H$. Successive
integration starts the three coordinates at Taylor orders $3,2,1$.
For $p>3$ and $n\geq k$, where $k=1,2,3$, one has
$n-v_p(n!)\geq k$. Thus the three coordinate valuations are at least
\[
 3(s-1)+3,\qquad3(s-1)+2,\qquad3(s-1)+1.
\]
Multiplication by the three columns of $B_\sigma$ raises all three
error components to $p^{3s}$. The displayed change-of-lift identity
therefore gives the desired congruence for $B_\phi$.
This proves \eqref{I:ap:jet} on $A_H$. Both sides belong to $A^3$;
the injection $A/p\hookrightarrow A_H/p$ makes the localization
$p$-saturated, and yields \eqref{I:ap:jet} on $A$.

The inputs from crystalline theory in this argument are its integral
Frobenius and Verschiebung on a smooth elliptic family, horizontality,
and the Taylor change-of-lift identity. The normalizations, the unit
change of basis, and the extra three-digit precision are established
explicitly above.

\subsection{The supersingular congruence at every cutoff}

The CM hypothesis is unnecessary on the supersingular side. The explicit
level-six family supplies a stronger statement.

\begin{theorem}[Supersingular Ap\'ery residue disks]\label{I:thm:aperydisks}
Let $p>3$ and $t_0\in\mathbb Z_p$ satisfy
\[
 t_0(1-34t_0+t_0^2)\in\mathbb Z_p^\times.
\]
Choose a root $h_0$ of
\[
 72t_0h^2+(17t_0-1)h+t_0=0
\]
in the ring $\mathcal O$ of an unramified extension of degree at most two.
Assume the reduction of the explicit elliptic curve $E_{h_0}$ in
\eqref{I:ap:elliptic} is supersingular. Then, for every positive integer
$N$ divisible by $p^2$,
\begin{equation}\label{I:ap:ssdisks}
 P_N(t_0)-p^2P_{N/p^2}(t_0)
       \in p^{3v_p(N)-1}\mathbb Z_p.
\end{equation}
\end{theorem}
The integral Frobenius--Taylor identity \eqref{I:ap:jet} was proved in the
preceding section. No CM condition on
$t_0$ occurs in the argument below.

\begin{proof}
The discriminant of the displayed quadratic is $1-34t_0+t_0^2$, a unit.
Its leading and constant coefficients are units, since $p>3$ and $t_0$
is a unit. Thus its two roots are units in an unramified extension of
degree at most two. The identity $D(h_0)=h_0/t_0$ makes $D(h_0)$ a unit
as well. Consequently both the elliptic curve and its $D$-twist have
good reduction.

The untwisted family has a rational point of exact order six:
\[
 R=\left(h(1+9h),\frac{hD}{2}\right),\qquad
 2R=\left(h^2,\frac{h^2(1+8h)}2\right),\qquad 3R=(0,0).
\]
These formulas follow directly from the addition law. The assumptions
that $hD$ is a unit and $p>3$ preserve the exact order after reduction.

First suppose $h_0\in\mathbb Z_p$. Its supersingular reduction is
defined over $\mathbb F_p$. The Frobenius trace is divisible by $p$ and
has absolute value at most $2\sqrt p<p$ for $p>3$, hence it is zero.
This holds for its quadratic twist as well. Choose
\[
 \phi(h)=h_0+(h-h_0)^p.
\]
It is an integral Frobenius lift fixing $h_0$. On the rank-two fiber,
$M^2=-pI$, so on the jet module $B^2=p^2I$. The first row of $M$ is
divisible by $p$ by the Hodge filtration, hence the first row of
$V=pM^{-1}=-M$ is divisible by $p$, and that of $B$ by $p^2$.

Next suppose $h_0$ has residue degree two. Write $\iota$ for the
nontrivial automorphism of $\mathcal O/\mathbb Z_p$. Its two roots are
$h_0,h_0^\iota$. We may choose an integral Frobenius lift that exchanges
them exactly: divide $h^p$ by the monic quadratic minimal polynomial,
let $r_p(h)$ be the remainder, let $r(h)$ be the affine map exchanging
the roots, and put
\[
 \phi(h)=h^p-r_p(h)+r(h).
\]
The irreducibility of the quadratic modulo $p$ gives
$r_p\equiv r\pmod p$, so this is indeed a Frobenius lift.

The reduction of $E_{h_0}$ over $\mathbb F_{p^2}$ is supersingular.
Its Frobenius trace is $cp$ with
$c\in\{-2,-1,0,1,2\}$: supersingularity gives divisibility by $p$,
and the Weil bound gives this range. Its rational point of order six
implies
\[
 6\mid\#E_{h_0}(\mathbb F_{p^2})=p^2+1-cp,
 \qquad cp\equiv2\pmod6.
\]
As $p\equiv\pm1\pmod6$, the only possible value of $c$ in the stated
range is $c=2$ when $p\equiv1\pmod6$, or $c=-2$ when
$p\equiv-1\pmod6$. The $p^2$-Frobenius is therefore scalar $\pm p$.
Indeed its characteristic polynomial is $(T\mp p)^2$. If an elliptic
endomorphism is nonzero, it is an isogeny and cannot be nilpotent;
therefore the relation $(\pi\mp[p])^2=0$ forces $\pi=\pm[p]$ exactly.
A quadratic twist over $\mathbb F_{p^2}$ changes at most the sign, so
the same scalar assertion holds for $E_{h_0}^D$.

Let $M_0,M_1$ be the rank-two matrices at $h_0,h_0^\iota$ with compatible
bases. Their two-step product is $\epsilon pI$ for $\epsilon=\pm1$.
Hence $V_0=\epsilon M_1$ and $V_1=\epsilon M_0$. Since each $M_i$ has
first row divisible by $p$, so does each $V_i$. Consequently the jet
matrices satisfy
\[
 B_0B_1=p^2I,\qquad\text{the first row of }B_i\text{ belongs to }p^2\mathcal O^3.
\]
The changes from symmetric-square coordinates to jets cancel in the
product. Their first row is $(1,0,0)$, so they preserve the indicated
first-row divisibility.

Write $N=mp^s$ with $(m,p)=1$ and $s\geq2$. The trace argument applies
to every cutoff $mp^j$, and the residual identity is
$L P_{mp^j}=(mp^j)^3 t^{mp^j}(-A_{mp^j}+tA_{mp^j-1})$.
The proof of \eqref{I:ap:jet} therefore gives the same identity with
$Y_j$ formed from $P_{mp^j}(t(h))$, with error in $p^{3j}$.
Apply it twice at the two points (the same point in residue degree one).
Taking the first coordinate gives
\[
 P_{mp^s}(t_0)-p^2P_{mp^{s-2}}(t_0)
       \in p^{3s}\mathcal O+p^2p^{3(s-1)}\mathcal O
       =p^{3s-1}\mathcal O.
\]
Both polynomials have values in $\mathbb Z_p$, and the extension is
unramified. Intersecting with $\mathbb Z_p$ proves
\eqref{I:ap:ssdisks}. The argument is unchanged for every lift of the
same allowed residue class, so it covers entire supersingular residue
disks.
\end{proof}

\paragraph{First level and valuation consequence.}
The same first-row divisibility and the $s=1$ jet identity give
\[
 P_{mp}(t_0)\in p^2\mathbb Z_p\qquad(m\geq1).
\]
Combining this with \eqref{I:ap:ssdisks} shows that for every $p\mid N$,
\[
 v_p(P_N(t_0))\geq 2\left\lceil\frac{v_p(N)}2\right\rceil.
\]
In the usual prime-power notation, for all $m\geq1$ and $s\geq2$,
\[
 P_{mp^s}(t_0)-p^2P_{mp^{s-2}}(t_0)
 \in p^{3(s+v_p(m))-1}\mathbb Z_p.
\]
The extra $3v_p(m)$ follows by writing $m=p^a m_0$ with $(m_0,p)=1$
and applying the theorem at level $s+a$.
The exponent cannot uniformly be increased from $3s-1$ to $3s$:
the exact finite calculation at $p=7$, $t_0=1$, $m=1$, $s=2$ gives
\[
 P_{49}(1)-49\equiv 2\cdot7^5\pmod{7^6}.
\]

\subsection{The two rational CM specializations}

For $t_0=-1$, take $h_0=-1/6$. Then $D(h_0)=1/6$ and the curve
$E_{h_0}$ is isomorphic over $\mathbb Z[1/6]$ to $y^2=x^3+1$.
If $p\equiv2\pmod3$, the map $x\mapsto x^3$ permutes $\mathbb F_p$,
so the trace of Frobenius is zero. Thus the theorem gives
\[
 \sum_{k<mp^s}(-1)^kA_k-p^2\sum_{k<mp^{s-2}}(-1)^kA_k
 \in p^{3(s+v_p(m))-1}\mathbb Z_p
\]
for every such $p>3$, $m\geq1$, $s\geq2$.

For $t_0=1$, the cover points are
\[
 72h^2+16h+1=0,\qquad h_\pm=\frac{-4\pm\sqrt{-2}}{36}.
\]
At either root, write $a_2=(1+12h+24h^2)/4$ and
$b=h^3(1+9h)$. Direct reduction modulo the quadratic gives
$a_2^2=8b$. Hence $E_h$ is the quadratic twist with parameter $-a_2/4$
of
\[
 E_0:\quad y^2=x^3-4x^2+2x,\qquad\Delta(E_0)=512.
\]
Moreover $a_2=(1+10h)/6$ and $\operatorname{Norm}(a_2)=1/216$, so
this twist parameter is a unit at every $p>3$.

The quotient of $E_0$ by $(0,0)$ is the degree-two isogenous curve
$y^2=x^3+8x^2+8x$, its $-2$ quadratic twist. Over
$\mathbb Q(\sqrt{-2})$ this gives an endomorphism $J$ with differential
multiplier $\sqrt{-2}$, and hence $J^2=[-2]$. Its Galois conjugate is
$-J$. At a prime with $(-2/p)=-1$, Frobenius anticommutes with $J$;
therefore it has trace zero and $E_0$ is supersingular. The unit twists
remain supersingular. This proves
\[
 \sum_{k<mp^s}A_k-p^2\sum_{k<mp^{s-2}}A_k
 \in p^{3(s+v_p(m))-1}\mathbb Z_p
\]
for all $p\equiv5,7\pmod8$, $m\geq1$, and $s\geq2$.
The argument includes $p=5,7$: it uses the characteristic-zero explicit
unit twist, so it does not incorrectly exclude primes at which the
reduction of $j=8000$ becomes $0$ or $1728$.

\subsection{A finite supersingularity check}

The theorem's hypothesis can be tested directly from its displayed
elliptic curve, without an unproved identification with a truncated
Ap\'ery Hasse polynomial. Reduce $h_0$ to $\bar h$ in $\mathbb F_{p^2}$
(or $\mathbb F_p$ when possible), and form
\[
 \mathcal H_p(\bar h)=
 [X^{p-1}]\left(X^3+\frac{1+12\bar h+24\bar h^2}{4}X^2
                  +\bar h^3(1+9\bar h)X\right)^{(p-1)/2}.
\]
The standard Cartier--Manin formula for a smooth elliptic curve says
that it is supersingular exactly when $\mathcal H_p(\bar h)=0$.
The extra $D$-twist multiplies this coefficient by the nonzero factor
$D(\bar h)^{(p-1)/2}$ and does not change its vanishing.

\subsection{Sources and scope}

The modular parametrizations used above are given in Beukers,
\emph{Supercongruences using modular forms}, arXiv:2403.03301v3,
Appendix Sections 8.10 and 9.5 \cite{III:Beukers2026}. Its Corollary 1.20 is a first-level
ordinary Ap\'ery sum congruence. Sun--Ye, arXiv:2408.09776v4, Theorem 1.29
and Section 9 prove the first-level ordinary alternating congruence
and the two CM-value identities \cite{I:SunYe}. Bloch--Vanhove, arXiv:1309.5865,
Section 4 gives the equivalent sunset elliptic pencil \cite{I:BlochVanhove}.
The present supersingular result covers all positive cutoffs divisible
by $p^2$ throughout the stated smooth supersingular disks. It excludes
$p=2,3$, $t_0\equiv0$, and the two branch values modulo $p$; no assertion
about these excluded loci is made here.


\paragraph{Precise limits of the statements.}
The supersingular theorem is relative to the open set
$t_0(1-34t_0+t_0^2)\in\mathbb Z_p^\times$ and the explicit elliptic
supersingularity condition. It makes no assertion at $p=3$ or at the
excluded singular reductions. The two rational specializations are
corollaries of the whole-disk theorem, rather than the only parameter
values covered. The modulus $3s-1$ cannot uniformly be replaced by
$3s$, as the displayed exact $p=7$ example shows.

\appendix\paperanchors{I.app}

\section{Reproducible finite checks}

These checks are supplementary diagnostics; no finite computation is
used in the proof.  A block-check script tested every odd prime
$p\le43$, every $2\le d\le2p+2$, both signs satisfying $\chi=1$, and
$m=1,2,3$.  It found no failures in $2583$ instances of
\eqref{I:eq:block}, including nonunit cases.
An independently run check tested ordinary cases with
$2\le d\le12$ and both signs for $p=3,5,7$ through level $s=3$, and
for $p=11,13,17,19,29,37$ through level $s=2$.
All $178$ level-recursion instances passed.

The following complete Python 3 listing reproduces both counts and
checks.  It computes central binomial coefficients by exact integer
recurrence before reducing modulo $p^2$.  It requires no external
package, data file, or approximation to the unit root.

\begin{lstlisting}
def primes(bound):
    return [p for p in range(3, bound + 1, 2)
            if all(p % j for j in range(3, int(p**0.5) + 1, 2))]

def coefficients(length, modulus):
    b, inv4, q, out = 1, pow(4, -1, modulus), 1, []
    for k in range(length):
        out.append((b % modulus) * q % modulus)
        b = b * 2 * (2*k + 1) // (k + 1)
        q = q * inv4 % modulus
    return out

def terms(a, d, eps, modulus):
    return [pow(x, d, modulus) * (eps if k % 2 else 1) % modulus
            for k, x in enumerate(a)]

block_checks = 0
for p in primes(43):
    modulus, N = p*p, (p-1)//2
    a = coefficients(3*p, modulus)
    for d in range(2, 2*p + 3):
        for eps in (-1, 1):
            if (-1)**(d*N) * eps**N != 1:
                continue
            c = terms(a, d, eps, modulus)
            fp = sum(c[:p]) % modulus
            for m in (1, 2, 3):
                assert (sum(c[:m*p]) - fp*sum(c[:m])) % modulus == 0
                block_checks += 1

level_checks = 0
for p in (3, 5, 7, 11, 13, 17, 19, 29, 37):
    modulus = p*p
    max_level = 3 if p <= 7 else 2
    a = coefficients(p**max_level, modulus)
    for d in range(2, 13):
        for eps in (-1, 1):
            c = terms(a, d, eps, modulus)
            values = [sum(c[:p**s]) % modulus
                      for s in range(1, max_level + 1)]
            if values[0] % p == 0:
                continue
            for s in range(1, len(values)):
                assert (values[s] - values[0]*values[s-1]) % modulus == 0
                level_checks += 1

assert block_checks == 2583
assert level_checks == 178
print("Block checks:", block_checks, "; level checks:", level_checks)
\end{lstlisting}


\section{Diagnostics for the level-dependent bound}

An exact modular calculation tested
\[
 F_{p^{s+1}}(\eps)F_{p^{s-1}}(\eps)-F_{p^s}(\eps)^2
                   \equiv0\pmod{p^{\nu_d(s)}}
\]
in ordinary cases with $2\le d\le12$ and both signs.
The tested primes and maximum values of $s$ were
$(3,5),(5,4),(7,3),(11,2),(13,2),(17,2)$, respectively.
All $256$ instances passed.  These finite checks are supplementary and
are not used in any proof above.  They do not test optimality of the
exponent.
Using the \texttt{coefficients} and \texttt{terms} functions in the
verification appendix, the additional check is reproduced by:
\begin{lstlisting}
checks = 0
for p, maxs in ((3,5), (5,4), (7,3), (11,2), (13,2), (17,2)):
    modulus = p**(2*maxs + 1)
    a = coefficients(p**(maxs + 1), modulus)
    for d in range(2, 13):
        for eps in (-1, 1):
            c = terms(a, d, eps, modulus)
            f = [1] + [sum(c[:p**j]) % modulus
                       for j in range(1, maxs + 2)]
            if f[1] % p == 0:
                continue
            for s in range(1, maxs + 1):
                exponent = min(2*s, s+d-1)
                assert (f[s+1]*f[s-1] - f[s]**2) % p**exponent == 0
                checks += 1
assert checks == 256
print("Level-dependent cross-difference checks:", checks)
\end{lstlisting}


\section{Exact checks on supersingular Ap\'ery disks}

The following standalone Python~3 program selects the residue classes
using the Cartier--Manin coefficient of the displayed elliptic curve
over $\mathbb F_{p^2}$. It does not select them by assuming the Ap\'ery
congruence being tested. It checks all the selected smooth residue
classes for the listed primes through $43$, at lifts $t$ and $t+p$,
and also higher levels at several lifts. All $94$ exact checks pass.
The finite computation is a diagnostic for the theorem, not its proof.

\begin{lstlisting}[basicstyle=\ttfamily\fontsize{7.7}{8.7}\selectfont]
#!/usr/bin/env python3
"""Exact diagnostics for the supersingular Apery-disk theorem.

Supersingularity is tested on the displayed elliptic curve over F_(p^2),
not inferred from the Apery truncation being tested. No floating point.
"""
from math import comb


def add(x, y, p):
    return ((x[0]+y[0]) % p, (x[1]+y[1]) % p)


def scale(x, c, p):
    return (x[0]*c % p, x[1]*c % p)


def mul(x, y, p, d):
    return ((x[0]*y[0]+d*x[1]*y[1]) % p,
            (x[0]*y[1]+x[1]*y[0]) % p)


def power(x, k, p, d):
    ans = (1, 0)
    while k:
        if k & 1:
            ans = mul(ans, x, p, d)
        x = mul(x, x, p, d)
        k >>= 1
    return ans


def ss_residues(p):
    d = next(i for i in range(2, p) if pow(i, (p-1)//2, p) == p-1)
    roots = []
    for t in range(1, p):
        disc = (1-34*t+t*t) % p
        if not disc:
            continue
        if pow(disc, (p-1)//2, p) == 1:
            root = (next(i for i in range(p) if i*i % p == disc), 0)
        else:
            ratio = disc * pow(d, -1, p) % p
            root = (0, next(i for i in range(p) if i*i % p == ratio))
        h = scale(add((1-17*t, 0), root, p), pow(144*t, -1, p), p)
        h2 = mul(h, h, p, d)
        a = scale(add(add((1, 0), scale(h, 12, p), p), scale(h2, 24, p), p),
                  pow(4, -1, p), p)
        b = mul(power(h, 3, p, d), add((1, 0), scale(h, 9, p), p), p, d)
        m = (p-1)//2
        hasse = (0, 0)
        for j in range(m//2+1):
            term = mul(power(a, m-2*j, p, d), power(b, j, p, d), p, d)
            hasse = add(hasse, scale(term, comb(m, j)*comb(m-j, j), p), p)
        if hasse == (0, 0):
            roots.append(t)
    return roots


def apery_coefficients(length, modulus):
    last, a = 0, 1
    result = []
    for n in range(length):
        result.append(a % modulus)
        numerator = (2*n+1)*(17*n*(n+1)+5)*a-n**3*last
        denominator = (n+1)**3
        assert numerator % denominator == 0
        last, a = a, numerator//denominator
    return result


def evaluate(coefficients, t, modulus):
    value = 0
    for c in reversed(coefficients):
        value = (value*t+c) % modulus
    return value


def valuation_residue(value, p, precision):
    value %= p**precision
    if value == 0:
        return f">={precision}"
    v = 0
    while value % p == 0:
        value //= p
        v += 1
    return v


def main():
    rows = []
    precision = 10
    for p in (5, 7, 11, 13, 17, 19, 23, 29, 31, 37, 41, 43):
        modulus = p**precision
        coeffs = apery_coefficients(p*p, modulus)
        residues = ss_residues(p)
        for residue in residues:
            for t in (residue, residue+p):
                defect = evaluate(coeffs, t, modulus)-p*p
                val = valuation_residue(defect, p, precision)
                assert isinstance(val, str) or val >= 5
                rows.append(dict(p=p, t=t, s=2, m=1, defect_valuation=val))
        print(p, residues, flush=True)
    for p, t, depth in ((5, 6, 4), (5, 11, 4), (5, 4, 4),
                        (5, 9, 4), (7, 8, 3), (7, 15, 3),
                        (11, 10, 3), (11, 21, 3)):
        assert t % p in ss_residues(p)
        precision = 3*depth+4
        modulus = p**precision
        coeffs = apery_coefficients(p**depth, modulus)
        vals = {s:evaluate(coeffs[:p**s], t, modulus) for s in range(depth+1)}
        for s in range(2, depth+1):
            val = valuation_residue(vals[s]-p*p*vals[s-2], p, precision)
            assert isinstance(val, str) or val >= 3*s-1
            rows.append(dict(p=p, t=t, s=s, m=1, defect_valuation=val))
    assert len(rows) == 94
    print(f'{len(rows)} exact supersingular-disk checks passed')


if __name__ == '__main__':
    main()

\end{lstlisting}




\clearpage\endgroup

% Source: Beukers_Delaygue_All_Levels_Lifted_Reflection.tex
\begingroup

\openpaper{II}{All-Level Beukers--Delaygue Supercongruences via Lifted Reflection}{All-level lifted reflection}{10 November 2026}{{\small Research manuscript prepared for Keio\par}\smallskip }



\begin{abstract}
Let $F_L(\epsilon)=\sum_{n<L}((1/2)_n/n!)^d\epsilon^n$, where
$d\ge2$ and $\epsilon=\pm1$. For every odd prime $p$ for which
$F_p(\epsilon)$ is a $p$-adic unit, we prove
\[
 F_{mp^s}(\epsilon)-f_p(\epsilon)F_{mp^{s-1}}(\epsilon)
                  \in p^{2s}\mathbb Z_p
 \qquad(m,s\ge1).
\]
Here $f_p(\epsilon)$ is the limiting quotient of successive prime-power
truncations; the proof also constructs this limit. At $m=1$ this is
the full statement of Beukers--Delaygue's Conjecture~1.5.
The proof treats both ranges $s<d$ and $s\ge d$ without a restriction
on the size of $p$. It identifies the carry contribution in the
analytic digit operator as an exact weighted finite difference.
A two-component integral lift then has rank at most one modulo $p^2$,
which gives a two-digit gain at every normalized iteration.
All analytic interpolation and operator identities are proved here;
no unpublished supercongruence theorem is assumed.
\end{abstract}

\noindent\fbox{\begin{minipage}{0.94\textwidth}
\textbf{Scope and research status.}
The theorem covers every level of Conjecture~1.5 under its original
unit hypothesis, including $p=3$. It also allows every positive
truncation multiplier $m$, without assuming $F_m(\epsilon)$ is a unit.
The proof has received internal independent checks. This manuscript
does not claim completed external peer review or certified literature
priority. The stronger special exponents and modular identifications
in Conjectures~1.3 and~1.4 are separate statements.
\end{minipage}}
\medskip

\section{Statement and proof strategy}
\label{II:sec:bdx}

Fix an odd prime $p$, an integer $d\ge2$, and $\epsilon\in\{1,-1\}$.
Write
\[
 \alpha_n=\frac{(1/2)_n}{n!}=4^{-n}\binom{2n}{n},\qquad
 F_L(\epsilon)=\sum_{n=0}^{L-1}\alpha_n^d\epsilon^n,
 \qquad N=\frac{p-1}{2}.
\]
All coefficients belong to $\mathbb Z_p$. The following theorem proves
the mathematical statement of Beukers--Delaygue's Conjecture~1.5 \cite{I:BD}, with
the additional allowance of arbitrary truncation multipliers.

\begin{theorem}[All levels and all lengths]\label{II:bdx:main}
Assume $F_p(\epsilon)\in\mathbb Z_p^\times$.
The limit
\[
 f_p(\epsilon)=\lim_{s\to\infty}
       \frac{F_{p^s}(\epsilon)}{F_{p^{s-1}}(\epsilon)}
\]
exists in $\mathbb Z_p^\times$. For all integers $m,s\ge1$,
\begin{equation}\label{II:bdx:alllength}
 F_{mp^s}(\epsilon)-f_p(\epsilon)F_{mp^{s-1}}(\epsilon)
                      \in p^{2s}\mathbb Z_p.
\end{equation}
There is no unit condition on $F_m(\epsilon)$.
Equivalently, for every positive integer $K$ divisible by $p$,
\begin{equation}\label{II:bdx:allcutoff}
 F_K(\epsilon)-f_p(\epsilon)F_{K/p}(\epsilon)
                      \in p^{2v_p(K)}\mathbb Z_p.
\end{equation}
The theorem includes $p=3$ and imposes no condition relating $p$ and $d$.
\end{theorem}

The proof replaces a direct parity estimate by an operator on a larger
integral space. An exact finite-difference term records the failure of
the full digit operator to commute with reflection. The larger space
retains that term, and its operator contracts every normalized
nonconstant component by $p^2$ at each iteration. No inverse of a
finite-difference operator, and no assertion that its image is closed,
is required.

\section{Integral analytic digit functions}

Let $\mathcal A=\mathbb Z_p\langle a\rangle$ be the algebra of restricted
power series, with its Gauss norm. Divisibility in this algebra always
means coefficientwise divisibility. Integral affine substitutions and
evaluations at points of $\mathbb Z_p$ have norm at most one.

\begin{lemma}\label{II:bdx:digits}
There are functions $g_t\in\mathcal A$, $0\le t<p$, satisfying
\[
 g_t(a)=\frac{\alpha_{pa+t}}{\alpha_a}
 \quad\hbox{for every nonnegative integer }a.
\]
For $0\le t\le N$, there is a $\beta_t\in\mathbb Z_p$ such that
\[
 g_t(a)\equiv\alpha_t(1+p\beta_t a)\pmod{p^2\mathcal A}.
\]
In particular its constant term is $\alpha_t$ and
$g_t\equiv\alpha_t\pmod{p\mathcal A}$. For $N<t<p$,
\begin{equation}\label{II:bdx:bt}
 g_t(a)=p(a+1/2)b_t(a),\qquad b_t\in\mathcal A.
\end{equation}
\end{lemma}
\begin{proof}
Cancellation of factors divisible by $p$ gives, initially at
nonnegative integer $a$,
\[
 \frac{g_0(a+1)}{g_0(a)}
 =C\frac{\prod_{0\le t<p,\ t\ne N}(1+pa/(t+1/2))}
             {\prod_{i=1}^{p-1}(1+pa/i)},\qquad
 C=\frac{\prod_{0\le t<p,\ t\ne N}(t+1/2)}{(p-1)!}.
\]
Both lists of constant unit factors reduce to
$\mathbb F_p^\times$, so $C\in1+p\mathbb Z_p$ and, for $j\ge1$,
\[
 s_j=\sum_{0\le t<p,\ t\ne N}(t+1/2)^{-j}
                         -\sum_{i=1}^{p-1}i^{-j}\in p\mathbb Z_p.
\]
Define the polynomial
\[
 P_j(a)=\sum_{v=0}^{j}
    \left\{\begin{matrix}j\\v\end{matrix}\right\}
     \frac{a(a-1)\cdots(a-v)}{v+1}.
\]
The braces are Stirling numbers of the second kind. Thus
$P_j(a)=\sum_{u=0}^{a-1}u^j$ at nonnegative integers, and every coefficient
of $P_j$ has valuation at least $-\lfloor\log_p(j+1)\rfloor$.
The series
\[
 L(a)=a\log C+
    \sum_{j\ge1}\frac{(-1)^{j-1}p^js_j}{j}P_j(a)
\]
converges in $\mathcal A$, since its $j$th nonconstant summand has
coefficient valuations at least
\[
 j+1-v_p(j)-\lfloor\log_p(j+1)\rfloor\ge2,
\]
and these bounds tend to infinity. The inequality is immediate for
$j=1,2,3$; for $j\ge4$ it follows from $p\ge3$ and
$3^{j-1}\ge j(j+1)$. In particular
$L\in p\mathcal A$ and $L-a\log C\in p^2\mathcal A$.
Set $g_0=\exp L$. Its difference equation and value $g_0(0)=1$
give the required integer values, and
$g_0\equiv1+a\log C\pmod{p^2\mathcal A}$.

For $1\le t<p$, set
\[
 g_t(a)=g_0(a)4^{-t}
     \frac{\prod_{j=1}^{2t}(2pa+j)}{\prod_{j=1}^{t}(pa+j)^2}.
\]
All denominator constants are units. This product proves both the
integer-value assertion and the affine congruence when $t\le N$.
When $t>N$, its numerator contains the factor
$2pa+p=2p(a+1/2)$; cancelling this factor gives
\eqref{II:bdx:bt} in $\mathcal A$.
\end{proof}

\section{Both exact reflections}

Let $r(a)=-a-1/2$ and $r_1(a)=-a-3/2$.

\begin{lemma}\label{II:bdx:reflections}
The digit functions satisfy the following exact identities in
$\mathcal A$:
\begin{align}
 g_{N-t}(r(a))&=(-1)^N g_t(a)
                  &&(0\le t\le N),\label{II:bdx:lowerreflection}\\
 b_{p+N-t}(r_1(a))&=(-1)^{N+1}b_t(a)
                  &&(N<t<p).\label{II:bdx:upperreflection}
\end{align}
\end{lemma}
\begin{proof}
Put $N_j=(p^j-1)/2$. For each fixed nonnegative integer $n$,
\[
 \frac{\alpha_{N_j-n}}{\alpha_{N_j}}
   =\prod_{k=0}^{n-1}\frac{N_j-k}{N_j-k-1/2}
                 \longrightarrow\alpha_n.
\]
Also
\[
 \frac{\alpha_{N_j}}{\alpha_{N_{j-1}}}
      =g_N(N_{j-1})\longrightarrow C_*:=g_N(-1/2).
\]
For fixed $a\ge0$ and $0\le t\le N$, use
$a_j=N_{j-1}-a$ and $n=pa+t$. The equality
$pa_j+N-t=N_j-n$ and the preceding limits give
\[
 g_{N-t}(r(a))=C_*g_t(a).
\]
Continuity extends this equality to $\mathbb Z_p$, hence to an
identity of restricted power series. Applying it twice gives
$C_*^2=1$. Since $g_N\equiv\alpha_N\equiv(-1)^N\pmod p$,
one has $C_*=(-1)^N$. This proves \eqref{II:bdx:lowerreflection}.

For $N<t<p$, let $t^*=p+N-t$, which also lies strictly between $N$
and $p$. Now use $a_j=N_{j-1}-a-1$, so that
$pa_j+t^*=N_j-(pa+t)$. The same finite-product limits give
\[
 g_{t^*}(r_1(a))
   =(-1)^N\frac{\alpha_{pa+t}}{\alpha_{a+1}}
   =(-1)^N\frac{a+1}{a+1/2}g_t(a)
 \qquad(a\ge0).
\]
Using \eqref{II:bdx:bt} on both sides cancels the displayed factor and
gives \eqref{II:bdx:upperreflection}. This last identity is between
elements of $\mathcal A$, so continuity proves it everywhere.
\end{proof}

Put $\chi=(-1)^{dN}\epsilon^N$. The ordinary assumption forces
\begin{equation}\label{II:bdx:chi}
 \chi=1.
\end{equation}
Indeed $\alpha_t\equiv(-1)^t\binom Nt\pmod p$ for $0\le t\le N$,
so $\alpha_{N-t}\equiv(-1)^N\alpha_t\pmod p$.
The other digits are divisible by $p$. Reflection of the lower-digit
sum therefore gives $F_p(\epsilon)\equiv\chi F_p(\epsilon)\pmod p$.
Since $p$ is odd and the sum is a unit, \eqref{II:bdx:chi} follows.

\section{The finite-difference identity and exact symmetrization}

Define bounded integral operators on $\mathcal A$ by
\[
 (Sf)(a)=\sum_{t=0}^{p-1}\epsilon^t g_t(a)^d f(pa+t),
 \qquad (Rf)(a)=f(-a-1/2),
\]
and
\begin{equation}\label{II:bdx:D}
 (Df)(a)=a^d f(a)-\epsilon(a+1/2)^d f(a+1),
 \qquad (Uf)(a)=g_0(a)^d f(pa).
\end{equation}
Let
\[
 (Bf)(a)=\sum_{t=N+1}^{p-1}\epsilon^t b_t(a)^d f(pa+t),
 \qquad (E_-f)(a)=f(a-1).
\]
The upper-digit part of $S$ is consequently
$H=p^d(a+1/2)^dB$; all these operators are bounded and integral.

\begin{lemma}[Two exact operator identities]\label{II:bdx:identities}
One has
\begin{align}
 SD&=p^dDU,\label{II:bdx:SD}\\
 RSR-S&=\epsilon p^dDE_-B.\label{II:bdx:commutator}
\end{align}
Consequently, with
\begin{equation}\label{II:bdx:TV}
 T=\frac{S+RSR}{2},\qquad V=-\frac\epsilon2 E_-B,
\end{equation}
the operator $T$ commutes with $R$ and
\begin{equation}\label{II:bdx:split}
 S=T+p^dDV.
\end{equation}
\end{lemma}
\begin{proof}
The summands of $SDf$ telescope by the recurrence
$\alpha_{n+1}/\alpha_n=(n+1/2)/(n+1)$. The two surviving terms are
\[
 p^d a^d g_0(a)^d f(pa)
 -\epsilon p^d(a+1/2)^d g_0(a+1)^d f(p(a+1)),
\]
which proves \eqref{II:bdx:SD}. This may first be checked at nonnegative
integers; the displayed functions all belong to $\mathcal A$, so the
identity holds there. No division by $a+1$ is needed in the final
identity.

The lower-digit operator commutes with $R$, by
\eqref{II:bdx:lowerreflection}, \eqref{II:bdx:chi}, and
$r(pa+t)=pr(a)+(N-t)$.
For the upper digits, \eqref{II:bdx:upperreflection} and
$pr_1(a)+(p+N-t)=r(pa+t)$ give
\[
 (Bf)(r_1(a))=\epsilon(-1)^d(BRf)(a).
\]
It follows that
\[
 (RHf)(a)=\epsilon p^d a^d(BRf)(a-1),
 \qquad (HRf)(a)=p^d(a+1/2)^d(BRf)(a).
\]
Hence $RS-SR=\epsilon p^dDE_-BR$.
Multiplication on the right by $R$, with $R^2=1$, proves
\eqref{II:bdx:commutator}. The remaining assertions follow immediately.
All divisions used in \eqref{II:bdx:TV} are by the unit $2$.
\end{proof}

The finite-difference operator has a useful interpretation. For
\[
 \mathcal L_m(f)=\sum_{a=0}^{m-1}\alpha_a^d\epsilon^a f(a),
\]
direct telescoping gives
\[
 \mathcal L_m(Df)=-m^d\alpha_m^d\epsilon^m f(m).
\]
In particular this boundary term is divisible by $p^{ds}$ when
$m=p^s$. The proof below retains $Df$ in an integral state space
instead of discarding it at any finite stage.

\section{An integral lift with two-digit contraction}

Let $\mathcal A^+=\{e\in\mathcal A:Re=e\}$, and set
\[
 \mathcal M=\mathcal A^+\oplus\mathcal A,\qquad
 q(e,b)=e+Db.
\]
Equip this complete $\mathbb Z_p$-module with the maximum Gauss norm.
Define
\begin{equation}\label{II:bdx:lift}
 \mathcal K(e,b)=\bigl(Te,\ p^dVe+p^dUb\bigr).
\end{equation}
It is a bounded integral endomorphism because $T$ preserves
$\mathcal A^+$. Equations \eqref{II:bdx:SD} and \eqref{II:bdx:split} give
the exact intertwining identity
\begin{equation}\label{II:bdx:intertwine}
 q\mathcal K=S q.
\end{equation}

Put $c=-1/4$, $v_0=(1,0)$, and $\ell(e,b)=e(c)$.
Then $\ell(v_0)=1$, and $\mathcal M=\mathbb Z_pv_0\oplus\ker\ell$.
The crucial improvement is
\begin{equation}\label{II:bdx:rankone}
 (1-v_0\ell)\mathcal K(\mathcal M)\subseteq p^2\ker\ell,
 \qquad \ell(\mathcal K v_0)\in\mathbb Z_p^\times.
\end{equation}
To see the first assertion, Lemma~\ref{II:bdx:digits} and
$f(pa+t)\equiv f(t)+pa f'(t)\pmod{p^2\mathcal A}$ show that $Sf$
is affine modulo $p^2$ for every $f\in\mathcal A$; the upper digits
vanish modulo $p^2$ because $d\ge2$.
Since $R$ is an integral affine substitution, $RSRf$ is also affine
modulo $p^2$. For $e\in\mathcal A^+$, $Te$ is even under $a\mapsto-a-1/2$.
An affine even function modulo $p^2$ is constant modulo $p^2$,
so $Te-Te(c)\in p^2\mathcal A^+$. The second component of
\eqref{II:bdx:lift} is divisible by $p^d$, hence by $p^2$.
Finally $\ell(\mathcal K v_0)\equiv F_p(\epsilon)\pmod p$,
which proves the unit assertion.

\begin{lemma}[Normalized iteration]\label{II:bdx:iteration}
There is a unit $\lambda\in\mathbb Z_p^\times$ such that
\begin{equation}\label{II:bdx:statecongruence}
 \mathcal K^s v_0-\lambda\mathcal K^{s-1}v_0
                    \in p^{2s}\mathcal M\qquad(s\ge1).
\end{equation}
\end{lemma}
\begin{proof}
On $\mathcal Q=v_0+p^2\ker\ell$ define
$\mathcal N(x)=\mathcal Kx/\ell(\mathcal Kx)$.
By \eqref{II:bdx:rankone}, the denominator is a unit and
$\mathcal N(\mathcal Q)\subseteq\mathcal Q$.
For $x,y\in\mathcal Q$ one has the exact identity
\[
 \mathcal N(x)-\mathcal N(y)
 =\frac{\mathcal K(x-y)-\ell(\mathcal K(x-y))\mathcal N(y)}
              {\ell(\mathcal Kx)}.
\]
If $x-y\in p^r\ker\ell$, then the numerator equals
\[
 (1-v_0\ell)\mathcal K(x-y)
 -\ell(\mathcal K(x-y))(\mathcal N(y)-v_0),
\]
which belongs to $p^{r+2}\ker\ell$ by \eqref{II:bdx:rankone}.
Thus normalized iteration gains two powers of $p$ at every step.

Set $h_0=v_0$ and $h_s=\mathcal N(h_{s-1})$.
Then $h_s-h_{s-1}\in p^{2s}\ker\ell$, so $h_s$ converges to
$h_\infty\in\mathcal Q$. Write
\[
 Z_s=\ell(\mathcal K^s v_0),\quad
 z_s=Z_s/Z_{s-1}=\ell(\mathcal K h_{s-1}),\quad
 \lambda=\ell(\mathcal K h_\infty).
\]
Inductively $Z_s$ and $z_s$ are units and
$h_s=\mathcal K^s v_0/Z_s$. Since
$h_\infty-h_{s-1}\in p^{2s}\mathcal M$, one has
$z_s-\lambda\in p^{2s}\mathbb Z_p$.
The identity
\[
 \mathcal K^s v_0-\lambda\mathcal K^{s-1}v_0
 =Z_{s-1}\bigl(z_s(h_s-h_{s-1})
                       +(z_s-\lambda)h_{s-1}\bigr)
\]
now proves \eqref{II:bdx:statecongruence}.
\end{proof}

\section{Return to the original truncation lengths}

Applying the bounded integral map $q$ to
\eqref{II:bdx:statecongruence}, and using $q(v_0)=1$ and
\eqref{II:bdx:intertwine}, gives the uniform function congruence
\begin{equation}\label{II:bdx:functioncongruence}
 S^s1-\lambda S^{s-1}1\in p^{2s}\mathcal A.
\end{equation}
The exact digit expansion gives, for every $m\ge1$ and $s\ge0$,
\begin{equation}\label{II:bdx:prefix}
 \mathcal L_m(S^s1)=F_{mp^s}(\epsilon).
\end{equation}
Indeed an iterated digit sequence represents a unique
$0\le b<p^s$; the product of the $g_t$ factors telescopes to
$\alpha_{p^sa+b}/\alpha_a$. Since $p$ is odd, the accompanying sign
after multiplication by $\epsilon^a$ is $\epsilon^{p^sa+b}$.
This is an exact rational identity and needs no unit assumption on
$\alpha_a$ or $F_m(\epsilon)$.

For $m=1$, $F_{p^s}(\epsilon)=(S^s1)(0)$.
Since $h_s\in v_0+p^2\ker\ell$ and $q$ is integral,
$q(h_s)(0)\in1+p^2\mathbb Z_p$. Hence these truncations are units,
and \eqref{II:bdx:functioncongruence} implies
\[
 \frac{F_{p^s}(\epsilon)}{F_{p^{s-1}}(\epsilon)}
                    \equiv\lambda\pmod{p^{2s}}.
\]
This proves existence of the stated limit and identifies it with
$\lambda=f_p(\epsilon)$.
Finally apply the bounded integral functional $\mathcal L_m$ to
\eqref{II:bdx:functioncongruence} to obtain \eqref{II:bdx:alllength}.
Writing $K=mp^{v_p(K)}$ with $p\nmid m$ yields
\eqref{II:bdx:allcutoff} and completes the proof of
Theorem~\ref{II:bdx:main}.

\paragraph{Why the earlier parity estimate lost precision.}
In the original function space the upper digits feed a reflection-odd
component of size $p^d$, whose naive estimate subsequently gains only
one digit per iteration. Identity \eqref{II:bdx:commutator} identifies
this contribution as an exact $D$ term. The lifted space stores its
primitive as the second coordinate; \eqref{II:bdx:SD} makes that coordinate
gain $d\ge2$ digits at every iteration. The argument therefore removes
the obstruction rather than assuming its value vanishes at the
evaluation point $a=0$.


\section{Consequences and precise scope}

Taking $m=1$ in Theorem~\ref{II:bdx:main} gives the complete statement of
Conjecture~1.5 of \cite{I:BD}. Taking $s=1$ gives Conjecture~1.2, since
$F_1(\epsilon)=1$. In particular, the range $s\ge d$ left open by
the earlier direct parity estimate is included in the same induction
as the smaller levels.

\begin{center}
\begin{tabular}{@{}p{0.28\textwidth}p{0.62\textwidth}@{}}
\toprule
Parameter or condition & Range in this manuscript \\
\midrule
Prime & Every odd prime, including $3$ \\
Number of parameters & Every integer $d\ge2$ \\
Sign & $\epsilon\in\{1,-1\}$ \\
Ordinary hypothesis & $F_p(\epsilon)\in\mathbb Z_p^\times$ \\
Level & Every $s\ge1$, including $s\ge d$ \\
Truncation multiplier & Every $m\ge1$; $F_m(\epsilon)$ may be a nonunit \\
Proved exponent & $2s$, or $2v_p(K)$ for a cutoff $K$ divisible by $p$ \\
\bottomrule
\end{tabular}
\end{center}

The unit hypothesis is exactly the one in the original conjecture.
The theorem makes no assertion for $p=2$, for other values of
$\epsilon$, or for nonunit $F_p(\epsilon)$.
Neither the larger special exponents in Conjecture~1.3 nor the
modular-coefficient identifications in Conjecture~1.4 follow from
Theorem~\ref{II:bdx:main} alone.

The earlier direct parity bound
$\min(2s,s+d-1)$ records the upper-digit contribution only by its size.
The operator identity \eqref{II:bdx:SD} records additional exact structure:
the primitive of every carry correction acquires $p^d$ at its next
iteration. Keeping that primitive as a separate coordinate removes
the artificial one-digit loss. This is an argument for every level,
not a finite verification followed by extrapolation.

\appendix\paperanchors{II.app}
\section{Reproducible exact diagnostics}

The proof does not depend on the following computations. They test
the two necessary finite identities
\[
 F_{p^{s+1}}F_{p^{s-1}}-F_{p^s}^{2}\equiv0\pmod{p^{2s}}
\]
and
\[
 F_{mp^s}F_{p^{s-1}}-F_{p^s}F_{mp^{s-1}}
                          \equiv0\pmod{p^{2s}},
 \qquad m=2,3,4,
\]
with a common sign $\epsilon$ in every factor. We use the prime and
maximum-level pairs
\[
 (p,s_{\max})=(3,8),\ (5,6),\ (7,4),\ (11,3),
\]
with $2\le d\le8$ and both signs whenever $F_p(\epsilon)$ is a unit.
All $656$ checks pass.
Of these checks, $236$ are in the previously uncovered range $s\ge d$.

The program uses exact modular arithmetic. In updating $\alpha_n$ it
separates powers of $p$ before inverting denominator units, so no
division by a nonunit modulo $p^k$ is performed. The recorded valuations
are lower bounds when a residue is zero at the working precision.

\begin{lstlisting}
from time import perf_counter

def coeffs(p, maxn, prec):
    mod=p**prec
    u,v=1,0
    out=[]
    for n in range(maxn):
        out.append(u*pow(p,v,mod)%mod if v<prec else 0)
        a,b=2*n+1,2*n+2
        while a%p==0:
            a//=p;v+=1
        while b%p==0:
            b//=p;v-=1
        u=u*a*pow(b,-1,mod)%mod
        assert v>=0
    return out

def vp(n,p,prec):
    if n==0:return prec
    ans=0
    while n%p==0:ans+=1;n//=p
    return ans

start=perf_counter()
checks=0
counts=[]
for p,maxs in ((3,8),(5,6),(7,4),(11,3)):
    prec=2*maxs+3
    mod=p**prec
    largest=max(p**(maxs+1),4*p**maxs)
    a=coeffs(p,largest,prec)
    for d in range(2,9):
        for eps in (-1,1):
            prefix=[0]
            for n,x in enumerate(a):
                prefix.append((prefix[-1]+pow(x,d,mod)*(eps if n%2 else 1))%mod)
            if prefix[p]%p==0:continue
            vals=[prefix[p**s] for s in range(maxs+2)]
            min_excess=prec
            for s in range(1,maxs+1):
                defect=(vals[s+1]*vals[s-1]-vals[s]**2)%mod
                assert defect%p**(2*s)==0,(p,d,eps,s,'cross',vp(defect,p,prec))
                checks+=1
                min_excess=min(min_excess,vp(defect,p,prec)-2*s)
                for m in (2,3,4):
                    defect=(prefix[m*p**s]*vals[s-1]-vals[s]*prefix[m*p**(s-1)])%mod
                    assert defect%p**(2*s)==0,(p,d,eps,s,m,vp(defect,p,prec))
                    checks+=1
            counts.append((p,d,eps,maxs,min_excess))
    print('prime',p,'done',flush=True)
print('all exact checks passed:',checks)
print('seconds:',round(perf_counter()-start,2))
print('metadata (p,d,eps,maxs,min consecutive excess):',counts)

assert checks == 656
previously_uncovered = sum(4*max(0, maxs-d+1)
                          for p,d,eps,maxs,excess in counts)
assert previously_uncovered == 236
print("Checks in the former s >= d gap:", previously_uncovered)

\end{lstlisting}



\clearpage\endgroup



% Source: AESZ4_Level54_Attractor_Supercongruence_Dossier.tex
\begingroup
\newcommand{\Z}{\mathbb Z}
\newcommand{\Q}{\mathbb Q}
\newcommand{\C}{\mathbb C}
\newcommand{\F}{\mathbb F}
\newcommand{\Zp}{\mathbb Z_p}
\newcommand{\Qp}{\mathbb Q_p}
\newcommand{\Gal}{\operatorname{Gal}}
\newcommand{\Frob}{\operatorname{Frob}}
\newcommand{\ord}{\operatorname{ord}}
\newcommand{\Sym}{\operatorname{Sym}}
\newcommand{\Fil}{\operatorname{Fil}}
\newcommand{\Span}{\operatorname{span}}
\newcommand{\AESZ}{\mathrm{AESZ}}
\newcommand{\PF}{\mathcal L}
\newcommand{\eps}{\varepsilon}
\newcommand{\vp}{v_p}

\openpaper{III}{The AESZ4 Attractor and Its Residual Monodromy}{The AESZ4 attractor and residual monodromy}{10 November 2026, revised for this edition}{{\small QEEDDD Verifiable-Mathematical-Reasoning-Corpus Project\par}\smallskip }
\allowdisplaybreaks\setlist{nosep}

\begin{abstract}
For the fourth-order hypergeometric system AESZ4, we compute the exact
Picard--Fuchs error of every truncation and classify the residual
monodromy of the standard companion lattice at every auxiliary prime.
The geometric action is irreducible modulo $\ell\ne3$. Modulo $3$
it is a single regular unipotent block, with a unique invariant flag.
A modular-polynomial calculation excludes a generic fiberwise
$2$-isogeny between the two elliptic factors. We record the ordinary
unit-root criterion that motivated the search for a rank-two summand,
including the integral comparison that this criterion requires.
Part~IV constructs the summand at the attractor and proves the
special-value congruence for every $p\ge5$.
\end{abstract}

\section{The two Frobenius comparisons}

Let $V_{z_*}$ denote the rank-four constituent of the third cohomology
selected by the AESZ4 family. Its middle cohomology in the Hesse
fiber-product model has dimension $40$. Two comparisons enter the
proof of Beukers's congruence. The first identifies a rank-two
summand of $V_{z_*}$, its holomorphic line, and its Frobenius
polynomial $X^2-a_p X+p^3$. The second connects the action of
Frobenius on this line to the finite hypergeometric sum modulo $p^3$.

Theorem~\ref{III:thm:residual} below controls the differential equation
at the cutoff. Theorem~\ref{III:thm:resmon} classifies the residual
monodromy of the varying family. At $z_*=-1/5832$ there are additional
symmetries of the special fiber. Part~IV, Appendix~\ref{IV:app:tetra},
uses them to construct the required summand, and its
Lemma~\ref{IV:lem:div} supplies the integral comparison.

\section{The conjecture, the family, and normalizations}

Define
\[
 F(t)={}_4F_3\!\left(
 \begin{matrix}\frac13,\frac13,\frac23,\frac23\\1,1,1\end{matrix};t\right)
 =\sum_{n\ge0}c_nt^n,
 \qquad
 c_n=\frac{(1/3)_n^2(2/3)_n^2}{(n!)^4}.
\]
For an integer $N\ge1$ write
\[
 P_N(t)=\sum_{n=0}^{N-1}c_nt^n.
\]
Because
\[
 \binom{-1/3}{k}^2\binom{-2/3}{k}^2
 =\frac{(1/3)_k^2(2/3)_k^2}{(k!)^4},
\]
Beukers's sum is exactly $P_p(-1/8)$.

\begin{conjecture}[Good-prime form of Beukers, Conjecture 1.28]\label{III:conj:beukers}
For every odd prime $p\neq3$,
\[
 P_p\!\left(-\frac18\right)\equiv a_p\pmod{p^3},
\]
where $a_p$ is the coefficient of $q^p$ in the rational newform
\[
 f_{54}(q)=q+2q^2+4q^4+3q^5+29q^7+8q^8+\cdots
\]
of weight $4$ and level $54$, LMFDB label $54.4.a.c$.
\end{conjecture}


The printed statement in \cite[Conjecture~1.28]{III:Beukers2026} has
$p\ne3$. At $p=2$ the sum is $161/162$, whereas $a_2=2$.
Their difference has $2$-adic valuation $-1$. The good-prime statement
above is proved in Part~IV, Theorem~\ref{IV:thm:main}.

The hypergeometric coordinate is related to the mirror-bicubic coordinate $z$ by
\[
 t=729z.
\]
Thus the attractor point
\[
 z_*=-\frac1{5832}
\]
corresponds to
\[
 t_*=-\frac18.
\]
The modular decomposition of the rank-four constituent uses the weight-$4$ form $54.4.a.c$ and the weight-$2$ form
$54.2.a.a$.


\section{Excellent lifts and specialization}

Coefficientwise congruences of Dwork type motivated the original
ordinary approach \cite{III:BV3,III:CompleteEdition}. For example, the
formal descent one might seek for the Hadamard square of the cubic
simplicial period has the shape
\[
 \frac{P_{mp^s}(t)}{F(t)}\equiv
 \frac{P_{mp^{s-1}}(\sigma(t))}{F(\sigma(t))}
 \pmod{p^{2s}} .
\]
Here $m,s\ge1$ and $\sigma$ denotes a prospective excellent
Frobenius lift. Using such an identity at $t_*$ would require an
analytic ring admitting that specialization, a comparison of $t_*$
with $\sigma(t_*)$, and an arithmetic interpretation of the
multiplier. At $m=s=1$ the displayed exponent is two,
whereas Beukers asks for three. Part~IV obtains the required
precision directly from factorial blocks and the rational lift
$z_*$ of the special fiber.

\section{The exact order-four cutoff residual}

Let $\theta=t\,d/dt$ and set
\[
 \PF=\theta^4-t\left(\theta+\frac13\right)^2
                  \left(\theta+\frac23\right)^2.
\]
Then $\PF F=0$.

\begin{theorem}[Exact cutoff residual]\label{III:thm:residual}
For every $N\ge1$,
\[
 \boxed{\PF P_N(t)=-N^4c_Nt^N.}
\]
In particular, for $N=p^s$ the residual is divisible by $p^{4s}$ in every coefficient ring in which
$c_Nt^N$ is integral.
\end{theorem}

\begin{proof}
The coefficient recurrence for $F$ is
\[
 n^4c_n=\left(n-\frac23\right)^2
        \left(n-\frac13\right)^2c_{n-1}\qquad(n\ge1).
\]
Applying $\PF$ to $P_N$ makes the coefficient of $t^n$ vanish for $0\le n<N$.
The only uncancelled term is produced by the shifted part at $n=N$, and equals
$-N^4c_Nt^N$.
\end{proof}


For $N=p^s$, the differential residual has four powers of $p$ before
any comparison with a Frobenius solution. Estimating the error of
such a comparison on the special fiber is a separate problem.

\section{An ordinary unit-root criterion}

In the rank-four Hodge structure at $z_*$ the decomposition has the
form
\[
 V_{z_*,\mathrm B}=\Lambda\oplus\Lambda^\perp,\qquad
 \Lambda_\C=V^{3,0}\oplus V^{0,3},\qquad
 (\Lambda^\perp)_\C=V^{2,1}\oplus V^{1,2}.
\]
Here $\mathrm B$ denotes the rational Betti realization, and
$V^{i,j}$ denotes its Hodge components. We first record what the
unit-root calculation gives under an explicit integral hypothesis.

\begin{definition}[Arithmetic splitting package]\label{III:def:package}
Fix an ordinary good prime $p\ge5$ with $p\nmid a_p$. Assume the following data.
\begin{enumerate}[label=(A\arabic*)]
\item an integral rank-two Frobenius-stable direct summand $W_p$ of the rank-four crystalline module at
$t_*$, containing the holomorphic period direction;
\item Hodge weights $\{0,3\}$ on $W_p$;
\item characteristic polynomial
\[
 \det(1-T\varphi\mid W_p)=1-a_pT+p^3T^2;
\]

\item an integral comparison of the truncated jet
$Y_p=(P_p,P_p',P_p'',P_p''')^{\mathsf t}$ at $t_*$
with the Frobenius unit eigenvalue $\alpha_p$ on $W_p$,
with first-coordinate error $P_p(t_*)-\alpha_p\in p^3\Zp$.

\end{enumerate}
\end{definition}

\begin{theorem}[Conditional splitting-to-supercongruence bridge]\label{III:thm:conditional}

Let $p\ge5$ and $p\nmid a_p$. If the hypotheses of
Definition~\ref{III:def:package} hold, then
\[
 P_p(t_*)\equiv a_p\pmod{p^3}.
\]
\end{theorem}

\begin{proof}
Modulo $p$ the polynomial $X^2-a_p X+p^3$ has the simple nonzero
root $a_p\bmod p$. Hensel's lemma gives its unit root $\alpha_p$,
and $a_p-\alpha_p=p^3/\alpha_p\in p^3\Zp$. Adding the
first-coordinate estimate in (A4) gives the assertion.
\end{proof}

\begin{remark}
The estimate in (A4) is an independent integral hypothesis. The
rational splitting and knowledge of its characteristic polynomial
identify the unit root at ordinary primes. They supply no estimate of
the truncation error. If $p\mid a_p$, the reduction of the polynomial
is $X^2$ and there is no unit root. Part~IV uses Cayley--Hamilton on
the holomorphic vector, so its result includes this case.
\end{remark}

\section{The exact arithmetic target}

The decomposition of the rank-four constituent identified in Part~IV,
Theorem~\ref{IV:thm:comparison}, is
\begin{equation}\label{III:eq:global-split}
 V_{z_*,\ell}^{\mathrm{ss}}
 \simeq
 \rho_{54.4.a.c,\ell}\oplus\rho_{54.2.a.a,\ell}(-1).
\end{equation}
Equivalently, at every good prime,
\begin{equation}\label{III:eq:euler-factor}
 \boxed{
 P_{3,p}(T)=
 (1-a_pT+p^3T^2)(1-pb_pT+p^3T^2).
 }
\end{equation}

Here $P_{3,p}$ is the Euler polynomial of $V_{z_*,\ell}$. Period
computations in \cite{III:BKSZ} identified the candidate modular
forms. A rational Hodge splitting alone does not establish the
Galois action on those summands.

Two approaches to the semisimplified representation are relevant.
One can construct a projector over $\Q$ and identify the
representations on its image and complement. Alternatively,
a proof of \eqref{III:eq:euler-factor} at all good primes would
identify the semisimplification by Chebotarev and Brauer--Nesbitt.
The latter comparison alone does not locate the holomorphic vector
in a specified de Rham summand. The tetrahedral projector in
Part~IV does locate this vector.

\section{Residual monodromy at every auxiliary prime}

The hypergeometric local system has standard companion matrices associated to
\[
 f(x)=(x^2+x+1)^2,
 \qquad
 g(x)=(x-1)^4.
\]
We use
\[
 A=\begin{pmatrix}
0&0&0&-1\\1&0&0&-2\\0&1&0&-3\\0&0&1&-2
\end{pmatrix},\qquad
 B=\begin{pmatrix}
0&0&0&-1\\1&0&0&4\\0&1&0&-6\\0&0&1&4
\end{pmatrix}.
\]
Put $e_4^T=(0,0,0,1)$ and $v=(0,-2,1,-2)^T$.  Then
\[
 A-B=3ve_4^T.
\]

\begin{theorem}[Complete residual-monodromy classification]\label{III:thm:resmon}
Let $\ell$ be a prime and let $V=\F_\ell^4$.
\begin{enumerate}[label=(\alph*)]
\item If $\ell\neq3$, the group generated by the reductions of $A$ and $B$ acts irreducibly on $V$.
\item If $\ell=3$, then $A\equiv B$, the common matrix is regular unipotent with characteristic and
minimal polynomial $(x-1)^4$, and its invariant subspaces are exactly the unique flag
\[
 0\subset\ker(A-I)\subset\ker(A-I)^2\subset\ker(A-I)^3\subset V.
\]
In particular there is no residual direct-sum decomposition $V=V_2\oplus V_2'$ preserved by the
geometric monodromy.
\end{enumerate}
\end{theorem}

\begin{proof}
Consider the observability matrix
\[
 \mathcal O=\begin{pmatrix}
 e_4^T\\e_4^TA\\e_4^TA^2\\e_4^TA^3
 \end{pmatrix}
 =\begin{pmatrix}
0&0&0&1\\0&0&1&-2\\0&1&-2&1\\1&-2&1&2
\end{pmatrix},
\qquad \det\mathcal O=1.
\]
Also consider the cyclic matrix
\[
 \mathcal C=(v,Av,A^2v,A^3v),
 \qquad \det\mathcal C=81.
\]
Let $W\subset V$ be invariant under both $A$ and $B$, and assume $\ell\neq3$.
If $e_4^T$ vanishes on $W$, then it vanishes on $A^iW$ for $0\le i\le3$; invertibility of
$\mathcal O$ gives $W=0$.  Otherwise choose $w\in W$ with $e_4^Tw\neq0$.  Since
$(A-B)w=3(e_4^Tw)v$, one has $v\in W$.  The determinant $81$ is nonzero modulo $\ell$, so
$v,Av,A^2v,A^3v$ span $V$ and $W=V$.

For $\ell=3$, one has $A=B$ and $A$ is the companion matrix of $(x-1)^4$.  Hence
$V$ is the cyclic $\F_3[x]$-module $\F_3[x]/(x-1)^4$.  Its submodules are the ideals
$(x-1)^j$, which form the displayed unique flag.
\end{proof}

\begin{corollary}[Uniqueness of the residual route]\label{III:cor:normalizer}
Suppose an integral arithmetic hypergeometric realization has geometric mod-$3$ monodromy generated by
the reductions of $A$ and $B$.  Since geometric monodromy is normal in arithmetic monodromy, the
arithmetic image normalizes it.  The unique flag in Theorem~\ref{III:thm:resmon}(b) is characteristic and is
therefore preserved by the arithmetic image.


For $\ell\neq3$, the standard companion lattice admits no rank-two
sub-local system after reduction. For $\ell=3$ its geometric
invariant subspaces form the stated flag.
\end{corollary}

This calculation completes the auxiliary-prime search on the specified
lattice. Lifting the flag to characteristic zero would require
additional information about extension classes. Specializing the family
can introduce new endomorphisms, as occurs for the tetrahedral
symmetries at $z_*$.

\section{The fiber-product and isogeny routes}

The Hadamard-square presentation suggests a fiber-product model built from the cubic elliptic family.
This makes low-degree fiberwise isogenies tempting.  The simplest versions fail.

\subsection{The Fricke-3 mismatch}
In the standard Hesse coordinate $\alpha$, the Fricke involution is
\[
 \alpha\longmapsto1-\alpha,
\]
whereas the involution induced by the proposed attractor fiber product is
\[
 \alpha\longmapsto-\frac1{8\alpha}.
\]
These are not the same rational map.  Thus the splitting is not explained by the naive identification of
the two elliptic factors through the standard Fricke $3$-isogeny.

\subsection{An exact degree-2 obstruction}
The Hesse $j$-function in this normalization is
\[
 j(\alpha)=\frac{27(1+8\alpha)^3}{\alpha(1-\alpha)^3}.
\]
Let $\Phi_2(X,Y)$ denote the classical level-$2$ modular polynomial.

\begin{proposition}[No generic degree-$2$ fiberwise isogeny]\label{III:prop:phi2}
The rational function
\[
 \Phi_2\!\left(j(\alpha),j\!\left(-\frac1{8\alpha}\right)\right)
\]
is nonzero.  More precisely, its numerator is
\[
\begin{split}
-19683&(8\alpha^2+20\alpha-1)^4
(64\alpha^4-112\alpha^3+60\alpha^2+14\alpha+1)\\
&\times Q_{12}(\alpha),
\end{split}
\]
where
\[
\begin{split}
Q_{12}(\alpha)={}&262144\alpha^{12}+189726720\alpha^{11}+69710610432\alpha^{10}
-232501817344\alpha^9\\
&+254561642496\alpha^8-48116259072\alpha^7-43846988736\alpha^6
+6014532384\alpha^5\\
&+3977525664\alpha^4+454105112\alpha^3+17019192\alpha^2-5790\alpha+1.
\end{split}
\]
The denominator is $\alpha^3(\alpha-1)^9(8\alpha+1)^9$.
\end{proposition}

\begin{proof}
Substitute the displayed $j$-functions into the classical polynomial
\[
\begin{split}
\Phi_2(X,Y)={}&X^3+Y^3-X^2Y^2+1488XY(X+Y)-162000(X^2+Y^2)\\
&+40773375XY+8748000000(X+Y)-157464000000000
\end{split}
\]
and factor over $\Z[\alpha]$.  Appendix~\ref{III:app:code} reproduces the exact computation.
\end{proof}

The obstruction concerns a generic degree-$2$ isogeny of the two elliptic fibers.
The correspondence in Part~IV acts simultaneously on the pencils and their common base.


\section{Residual hypotheses in a Frobenius comparison}

At a fixed fiber, a finite set of Frobenius traces can suffice after
the determinant, the residual representation, and its possible
deviations have been determined. A mod-$3$ flag for the generic
monodromy does not give those data at $z_*$. A naive Livn\'e test
based only on quadratic characters unramified outside $2,3$ therefore
leaves an implication to prove. Theorem~\ref{IV:thm:comparison} in
Part~IV instead compares the already constructed rank-two
representations at the auxiliary prime $2$. Its residual image and the
possible deviation fields specify the finite set of test primes.

\section{A criterion using a Hodge projector}

Let $e_{\mathrm H}\in\operatorname{End}_{\mathrm{HS}}(V_{z_*,\mathrm B})$
denote the projector onto $V^{3,0}\oplus V^{0,3}$ along its
orthogonal complement. Extending it by zero on the orthogonal
complement of $V_{z_*,\mathrm B}$ defines a Hodge class in
$H^6(X_{z_*}\times X_{z_*},\Q)(3)$.

\begin{proposition}[Compatible realizations of a projector]\label{III:prop:absolute}
Suppose $e_{\mathrm H}$ is realized by a motivated projector
defined over $\Q$, or by a descended absolute-Hodge projector with
compatible Galois and crystalline realizations. Then its rank-two
image is stable in the $\ell$-adic realizations and under Frobenius
in the good crystalline fibers. At an ordinary good prime, if these
summands have the stated modular polynomial and the integral
hypotheses of Definition~\ref{III:def:package} hold, then
$P_p(t_*)\equiv a_p\pmod{p^3}$.
\end{proposition}

\begin{proof}
Realization preserves the idempotent equation and its rank. Over
$\Q$, the realization of the projector commutes with the Galois
action and, at good primes, with Frobenius. The final assertion
uses the integral comparison in (A4) and
Theorem~\ref{III:thm:conditional}.
\end{proof}

This criterion records an alternative to finding a correspondence
explicitly. Appendix~\ref{IV:app:tetra} gives an algebraic
correspondence on the parabolic constituent.

\section{The finite-field hypergeometric route}

A direct proof of
\[
 P_{3,p}(T)=
 (1-a_pT+p^3T^2)(1-pb_pT+p^3T^2)
\]
would bypass the global projector.  Hypergeometric motives admit trace descriptions through Gauss and
Jacobi sums, and the special value $-1/8$ suggests cubic and quadratic transformations.  The search
found finite-field Clausen and Kummer identities near this value, but not an identity that factors the
entire degree-four Euler polynomial into the two specified modular quadratics.

The exact statement needed is stronger than a single trace identity.  One must control enough power
traces (or trace and determinant data on a rank-two summand) to identify the whole quartic factor at every
good prime.

\section{Exact numerical checks}

The verification program in Appendix~\ref{III:app:code} computes the truncation in exact rational arithmetic
and reduces it modulo $p^3$.  The first residues are
\[
\begin{array}{c|rrrrrrrrr}
p&5&7&11&13&17&19&23&29&31\\ \hline
P_p(-1/8)\bmod p^3&3&29&-57&20&-72&-106&174&-210&47.
\end{array}
\]
For $p=5,7$, these equal the displayed coefficients $a_5=3$ and $a_7=29$ of $54.4.a.c$.
The remaining entries are recorded as exact project checks; their identification with the corresponding
LMFDB coefficients is not used as a proof.

\section{Development of the arithmetic argument}


\begin{enumerate}[leftmargin=*,label=\textup{(\roman*)}]
\item The formal descent problem suggested isolating the special
fiber as the source of the extra power of $p$.
\item The Picard--Fuchs recurrence gave the exact residual
$-N^4c_Nt^N$. The ordinary rank-two polynomial separated comparison
with the truncated period from comparison with the Fourier coefficient.
\item Numerical period and Frobenius data suggested the weight-four
and twisted weight-two factors of \eqref{III:eq:euler-factor}.
\item The Fricke involution and the degree-two modular polynomial
eliminated two simple fiberwise explanations. The companion
matrices then gave the generic residual classification at every
auxiliary prime.
\item The simultaneous transformations of the Hesse pencils at
$z_*$ produced a rational projector. The four-prime comparison on the
rank-two summands identified their modular factors. Interpolation of
the shifted factorials in Part~IV completed the integral comparison
for every good prime.
\end{enumerate}

\section{From the projector to the truncation}

The two results in Part~IV can be used separately. Its
Theorem~\ref{IV:thm:comparison} gives \eqref{III:eq:global-split}
and \eqref{III:eq:euler-factor}. Its Lemma~\ref{IV:lem:div}
identifies the first entry of $p^3\Phi^{-1}$ on the rational lift
with $P_p(t_*)$ modulo $p^3$, and controls the integral lattice in
which the characteristic polynomial acts. Together they give
Theorem~\ref{IV:thm:main}. This route applies to ordinary and
nonordinary good primes alike.

\appendix\paperanchors{III.app}
\section{Exact verifier}\label{III:app:code}

The following program verifies the coefficient recurrence, the first truncation residues, the determinants
used in Theorem~\ref{III:thm:resmon}, and the exact modular-polynomial factorization of
Proposition~\ref{III:prop:phi2}.  It uses exact arithmetic throughout.

\begin{lstlisting}[style=python]
#!/usr/bin/env python3
"""Exact checks for the AESZ4 level-54 attractor dossier.

Standard-library checks are used where possible; SymPy is used only for exact
symbolic polynomial factorization and matrix determinants.
"""
from __future__ import annotations

import json
from fractions import Fraction
from math import factorial
from pathlib import Path

import sympy as sp

OUT = Path(__file__).with_name("verification_output.json")


def generalized_binomial(a: Fraction, k: int) -> Fraction:
    result = Fraction(1)
    for j in range(k):
        result *= a - j
    return result / factorial(k)


def truncated_sum_mod(p: int) -> int:
    modulus = p**3
    total = Fraction(0)
    for k in range(p):
        total += (
            generalized_binomial(Fraction(-1, 3), k) ** 2
            * generalized_binomial(Fraction(-2, 3), k) ** 2
            * Fraction(-1, 8) ** k
        )
    return (total.numerator * pow(total.denominator, -1, modulus)) % modulus


def picard_fuchs_check(max_n: int = 12) -> bool:
    # c_n=(1/3)_n^2(2/3)_n^2/(n!)^4 and
    # n^4 c_n=(n-2/3)^2(n-1/3)^2 c_{n-1}.
    c = [Fraction(1)]
    for n in range(1, max_n + 1):
        c.append(
            c[-1]
            * Fraction(3 * n - 2, 3) ** 2
            * Fraction(3 * n - 1, 3) ** 2
            / n**4
        )
    return all(
        n**4 * c[n]
        == Fraction(3 * n - 2, 3) ** 2
        * Fraction(3 * n - 1, 3) ** 2
        * c[n - 1]
        for n in range(1, max_n + 1)
    )


def residual_monodromy_data() -> dict[str, object]:
    A = sp.Matrix(
        [
            [0, 0, 0, -1],
            [1, 0, 0, -2],
            [0, 1, 0, -3],
            [0, 0, 1, -2],
        ]
    )
    B = sp.Matrix(
        [
            [0, 0, 0, -1],
            [1, 0, 0, 4],
            [0, 1, 0, -6],
            [0, 0, 1, 4],
        ]
    )
    e4 = sp.Matrix([[0, 0, 0, 1]])
    v = sp.Matrix([0, -2, 1, -2])
    observability = sp.Matrix.vstack(*[e4 * (A**i) for i in range(4)])
    cyclic = sp.Matrix.hstack(*[(A**i) * v for i in range(4)])
    assert A - B == 3 * v * e4
    assert observability.det() == 1
    assert cyclic.det() == 81
    x = sp.symbols("x")
    assert sp.factor(A.charpoly(x).as_expr()) == (x**2 + x + 1) ** 2
    assert sp.factor(B.charpoly(x).as_expr()) == (x - 1) ** 4
    return {
        "A_minus_B": [[int(z) for z in row] for row in (A - B).tolist()],
        "observability_det": int(observability.det()),
        "cyclic_det": int(cyclic.det()),
        "charpoly_A": str(sp.factor(A.charpoly(x).as_expr())),
        "charpoly_B": str(sp.factor(B.charpoly(x).as_expr())),
        "mod3_flag_dimensions": [1, 2, 3, 4],
    }


def phi2_factorization() -> dict[str, str]:
    a, X, Y = sp.symbols("a X Y")
    phi2 = (
        X**3
        + Y**3
        - X**2 * Y**2
        + 1488 * X * Y * (X + Y)
        - 162000 * (X**2 + Y**2)
        + 40773375 * X * Y
        + 8748000000 * (X + Y)
        - 157464000000000
    )
    j = lambda z: 27 * (1 + 8 * z) ** 3 / (z * (1 - z) ** 3)
    expression = sp.factor(
        sp.together(phi2.subs({X: j(a), Y: j(-sp.Rational(1, 8) / a)}))
    )
    numerator, denominator = sp.fraction(expression)
    assert numerator != 0
    return {
        "numerator": str(sp.factor(numerator)),
        "denominator": str(sp.factor(denominator)),
    }


def main() -> None:
    residues = {str(p): truncated_sum_mod(p) for p in [5, 7, 11, 13, 17, 19, 23, 29, 31]}
    expected_signed = {
        "5": 3,
        "7": 29,
        "11": -57,
        "13": 20,
        "17": -72,
        "19": -106,
        "23": 174,
        "29": -210,
        "31": 47,
    }
    for p_string, signed in expected_signed.items():
        p = int(p_string)
        assert residues[p_string] == signed % (p**3)

    report = {
        "picard_fuchs_coefficient_recurrence": picard_fuchs_check(),
        "truncated_sum_mod_p3": residues,
        "signed_representatives": expected_signed,
        "residual_monodromy": residual_monodromy_data(),
        "phi2_substitution": phi2_factorization(),
    }
    OUT.write_text(json.dumps(report, indent=2) + "\n", encoding="utf-8")
    print(json.dumps(report, indent=2))


if __name__ == "__main__":
    main()
\end{lstlisting}

\section{Verification output}
\begin{verbatim}
Picard--Fuchs recurrence: PASS
observability determinant: 1
cyclic determinant: 81
characteristic polynomial of A: (x^2+x+1)^2
characteristic polynomial of B: (x-1)^4
mod-3 invariant flag dimensions: 1,2,3,4
truncated residues (signed):
p=5: 3; p=7: 29; p=11: -57; p=13: 20; p=17: -72;
p=19: -106; p=23: 174; p=29: -210; p=31: 47.
Phi_2 substitution numerator: nonzero; exact factors given above.
\end{verbatim}

\section{Source and verification discipline}


The verifier checks the cutoff recurrence, the companion matrices, the
sample truncations, and the modular-polynomial calculation. The
arithmetic decomposition uses the geometric projector and the
Galois comparison in Part~IV. The verification range for the
good-prime congruence is recorded separately in that part.




\clearpage\endgroup

% Source: Beukers_128_Good_Primes.tex and unique finite identities from AESZ4_Tetrahedral_Projector.tex.
\begingroup
\newcommand{\Q}{\mathbf Q}\newcommand{\Z}{\mathbf Z}
\newcommand{\F}{\mathbf F}\newcommand{\A}{\mathcal A}
\newcommand{\Tr}{\operatorname{Tr}}\newcommand{\Gal}{\operatorname{Gal}}
\newcommand{\PP}{\mathbf P}\newcommand{\HH}{\mathcal H}
\newcommand{\Frob}{\operatorname{Frob}}
\openpaper{IV}{The Level-54 Attractor Supercongruence for Every Good Prime}{The level-54 good-prime theorem}{12 November 2026}{}
\begin{abstract}
We obtain the congruence of Beukers at the level-54 attractor for every
prime $p\ge5$, with no ordinarity hypothesis. Interpolating complete
blocks in the shifted factorial period gives an overconvergent inverse
Frobenius on the cyclic Picard--Fuchs frame. Its first coordinate is the
truncated period modulo $p^3$. Its four rows have divisibilities
$1,p,p^2,p^3$, and its determinant at a fixed fiber is $p^6$. These
facts turn the characteristic polynomial of the rank-two component
containing the holomorphic form into the required truncation
congruence. We include the construction and modular identification of
that component and an independent finite calculation. The prime $2$ in
the printed formulation fails by direct evaluation.
\end{abstract}

\section{The good-prime congruence}
Write $A_n=(3n)!/(n!)^3$, $c=-1/5832$, and $S_p=\sum_{n=0}^{p-1}A_n^2c^n$.
Let $a_p$ be the coefficient of $q^p$ in the normalized newform
$f$ labeled $54.4.\mathrm{a}.\mathrm{c}$. It has weight four, level
54, trivial character, and begins $q+2q^2+4q^4+3q^5+\cdots$.
\begin{theorem}\label{IV:thm:main}
For every prime $p\ge5$,
\[
  \sum_{k=0}^{p-1}\binom{-1/3}{k}^{\!2}\binom{-2/3}{k}^{\!2}
          \left(-\tfrac18\right)^k = S_p\equiv a_p\pmod{p^3}.
\]
\end{theorem}
The equality follows from $(1/3)_k(2/3)_k/(k!)^2=A_k/27^k$.
The theorem treats the good primes of Conjecture~1.28 of
Beukers~\cite{III:Beukers2026}. Its printed quantifier $p\ne3$ includes $2$,
where $S_2=161/162$, $a_2=2$, and $v_2(S_2-a_2)=-1$.

Let $V_z$ be the four-dimensional summand of B\"onisch's Hesse
fiber-product family \cite[\S5.1]{IV:bonisch}. Its ambient threefold has
40-dimensional $H^3$. At $z=c$, Appendix~\ref{IV:app:tetra} constructs
a rational rank-two projector whose image $W$ contains the
holomorphic form $\Omega_c$. Appendix~\ref{IV:app:modular} identifies
the trace and determinant of its Frobenius as $a_p,p^3$.
We use the convention that twisting by $(-1)$ multiplies good-prime
eigenvalues by $p$.

The passage from $W$ to the truncation occupies
Sections~\ref{IV:sec:shift}--\ref{IV:sec:fiber}. The matrix is built
over the whole smooth parameter locus; only its final
Cayley--Hamilton identity is imposed on the special fiber.
\section{A Frobenius operator from the shifted coefficients}\label{IV:sec:shift}
\subsection{Interpolation of the complete factorial blocks}
Write $\theta=z\,d/dz$ and use the rising factorial $(x)_n=\prod_{i=0}^{n-1}(x+i)$.
The differential equation in the factorial coordinate is
\begin{equation}\label{IV:eq:pf-shift}
 \theta^4 y=729z(\theta+\tfrac13)^2(\theta+\tfrac23)^2y .
\end{equation}
It is nonsingular on the base $\PP^1-\{0,1/729,\infty\}$.
All local estimates that follow include $p=5$.
Fix $p\ge5$ and put
\[
 P(T)=\prod_{\nu=1}^{p-1}(1+pT/\nu).
\]
Let $H(T)$ be the discrete primitive of $\log P(T)$, normalized by
$H(0)=0$. To construct it explicitly, write
\[
 \log P(T)=\sum_{m\ge1}d_m T^m,\quad
 d_m=(-1)^{m+1}\frac{p^m}{m}\sum_{\nu=1}^{p-1}\nu^{-m}.
\]
Replace each $T^m$ by
$\sum_{j=0}^m \genfrac{\{}{\}}{0pt}{}{m}{j}T^{\underline{j+1}}/(j+1)$, where the
braces are Stirling numbers of the second kind and the underline is
the falling factorial. Translation by one minus the identity sends
this polynomial to $T^m$. For $m=1,2$ the $d_m$ belong to
$p^3\Z_p$ by Wolstenholme and the power sum in $\F_p$.
For $m\ge3$ every coefficient of the replacement times $d_m$ has
valuation at least $m-v_p(m)-\lfloor\log_p(m+1)\rfloor\ge3$.
The lower bound grows linearly apart from logarithmic corrections and
the degree is at most $m+1$. Thus $H\in p^3\Z_p\langle T\rangle$
overconverges and has sufficiently small norm for exponentiation on
a disk strictly larger than the unit disk. More explicitly, for degree
$m+1$ each contribution to $H$ has valuation at least
$m-v_p(m)-\lfloor\log_p(m+1)\rfloor$, and the low degrees have
valuation at least three. Choosing a sufficiently small $\epsilon>0$
makes the Gauss norm of $H$ on $|T|\le p^\epsilon$ strictly smaller
than $p^{-1/(p-1)}$. This proves norm convergence of the exponential
in the corresponding Tate algebra. Define
\[
 b_0(T)=\exp(2[H(3T)-3H(T)]),\qquad
 b_k(T)=b_0(T)\left(\frac{(1+3pT)_{3k}}{(1+pT)_k^3}\right)^2
 \quad (1\le k<p).
\]
In particular $b_0(0)=1$, $b_k(0)=((3k)!/(k!)^3)^2$ and
$[T^d]b_k(T)\in p^{\min(d,3)}\Z_p$ for all $d\ge0$.
In this estimate the numerator is left unnormalized even where one
of its constants $1,\ldots,3k$ is divisible by $p$. All denominators
have unit constant. The factors are overconvergent on a common strict
neighborhood of the unit disk for any given $p$.

For $a_j(u)=((1+3u)_{3j}/(1+u)_j^3)^2$, splitting off factors
at multiples of $p$ gives the identity of germs at $u=0$
\[
 a_{pj+k}(pu)=a_j(u)\frac{b_k(j+u)}{b_0(u)} .
\]
Indeed the products of the nonmultiples of $p$ in the complete blocks
are products of $P(h+3u)$ and $P(h+u)$, respectively. The constant
$(p-1)!$ of each block cancels between numerator and denominator, and
summing the logarithms yields the quotient $b_0(j+u)/b_0(u)$.
The remaining $3k$ and $k$ factors yield $b_k(j+u)/b_0(j+u)$.
This proof works even for the indices $j$ with $p$-adic denominators
in their individual $u$-Taylor coefficients, since $a_j(u)$ has
been canceled before the block interpolation.

\subsection{The operator on the cyclic frame}
Use the cyclic vector $\omega=\Omega_z$ of the Picard--Fuchs equation, let
$\theta=z\,d/dz$, and $e_j=\theta^j\omega$, $0\le j\le3$.
Define an operator from the module at $z$ into its pullback at $z^p$
by putting $w=z^p$ and setting
\begin{equation}\label{IV:eq:shift-operator}
 C(e_j(z))=\left.\sum_{k=0}^{p-1} z^k
          b_k(\theta_w)(k+p\theta_w)^j\omega(w)\right|_{w=z^p}
 \qquad(0\le j\le3).
\end{equation}
Here $\theta$ inside the differential operator acts at the target
argument before substitution. The differential equation eliminates
$\theta^4\omega$ with coefficients in
$\Z_p[z,(1-729z)^{-1}]$. Each additional differentiation in this frame increases pole order and
numerator degree by at most a fixed constant, since the derivative of
a rational entry is added to its product with the companion matrix
of~\eqref{IV:eq:pf-shift}. Hence the preceding overconvergent series for each column
converges over the weak completion, by taking a sufficiently small
strict neighborhood of the good locus. The reductions do not lose
$p$-adic integrality. A column written in the ordered target frame has
its entries in $\Z_p,p\Z_p,p^2\Z_p,p^3\Z_p$,
respectively, at any good integral point. To see this directly, the
coefficient of $\theta_w^h$ in $b_k(\theta_w)(k+p\theta_w)^j$ has
valuation at least $\min(h,3)$. For $h=0,1,2$ the indicated monomial is
already a frame element, and for $h\ge3$ all terms after reduction are
divisible by $p^3$. In particular row $i$ of the matrix, where the
target frame element is $e_i(w)$, is divisible by $p^i$.
At precision $p^3$ its entry in row zero and column zero is
$\sum_{k=0}^{p-1}((3k)!/(k!)^3)^2 z^k$.

Write $G(z,u)=z^u\sum_{n\ge0}a_n(u)z^n$ modulo $u^4$, using formal
logarithms. Its four $u$-coefficients are a fundamental system of
solutions of~\eqref{IV:eq:pf-shift}: the indicial term is $u^4z^u$ and
the recurrence at each positive index is exactly the recurrence for
$a_n(u)$. Their leading logarithmic terms $1,\log z,(\log z)^2/2$ and
$(\log z)^3/6$ are linearly independent. Evaluating the target column against these four solution
jets is exactly the coefficient identity
\[
  \sum_{k=0}^{p-1}z^k
    [b_k(\theta)(k+p\theta)^jG(\,\cdot\,,u)](z^p)
  =b_0(u)\,\theta^j G(z,pu).
\]
All computations in this identity are coefficientwise at $z=0$ with
coefficients in $(\Q_p[u]/(u^4))[\log z]$. Running over the four solution jets gives an identity of fundamental
matrices. The matrix acting on the jets on the right is constant
in $z$, so differentiation proves horizontality, including the
relation at derivative order four. Its determinant is $p^6$, since
$u\mapsto pu$ acts by $1,p,p^2,p^3$ and multiplication by $b_0(u)$
has unit diagonal. The indicial initial conditions normalize the
$\theta$-Wronskian to one at $z=0$. Abel's identity applied to
\eqref{IV:eq:pf-shift} then gives the Wronskian $(1-729z)^{-2}$.
This gives $\det C=p^6(1-729z^p)^2/(1-729z)^2$ away from the deleted
disks. In particular after Taylor transport at a rational fiber the
determinant is $p^6$.
\section{Normalizing the geometric Frobenius}\label{IV:sec:normalization}
Write $\theta=z\,d/dz$ and $\lambda=729$. The algebraic connection of the
Hesse family on $B=\PP^1\setminus\{0,\lambda^{-1},\infty\}$ is the
rank-four hypergeometric connection used above. In particular the vector
$e_0$ in its cyclic frame is the normalized holomorphic three-form
$\Omega_z$. This normalization matters for the coefficient calculation.
B\"onisch \cite[\S5.1]{IV:bonisch} constructs $\Omega_z$ by residues of the two
cubic pencils, computes its period at $0$ as $\sum A_n^2z^n$, and
identifies $\Omega_z,\theta\Omega_z,\theta^2\Omega_z,\theta^3\Omega_z$ as
a cyclic frame. Its companion equation has no additional singular points.

Fix $p\ge5$. The Hesse surfaces and their projective fiber product
give a smooth projective family of threefolds over the smooth parameter
locus after reduction. To see good reduction in families,
one can first adjoin $\zeta_3$ and $z^{1/3}$ on this locus. The
base points of each Hesse pencil are nine distinct points and its
four singular fibers are triangles. The singular fibers of the two
pencils are disjoint precisely away from $\lambda z=1$ (and the
degenerate parameters $0,\infty$). Thus the total threefold in each
fiber is smooth. The vertical restriction kernels and the cyclic
average defining $V$ are functorial in these families and descend.
In particular relative cohomology supplies $V$ with a geometric
Frobenius $\Phi:\sigma^*V\to V$ for $\sigma(z)=z^p$. For clarity, the
smooth projective overconvergence theorem applies to the ambient
Gauss--Manin connection \cite[Corollary~6.4]{IV:lazda}. On a finite \'etale parameter cover splitting the components,
the restriction maps to the ruled components commute with this
Frobenius. Taking their kernel and cyclic invariants inside the
category of overconvergent isocrystals, followed by descent, gives
the same overconvergence for $V$. At a rational lift in a residue
disk, Taylor transport identifies $\Phi$ with the Frobenius on the
crystalline cohomology of that fiber in the de Rham realization.

\begin{lemma}\label{IV:lem:normal}
The horizontal map $C:V\longrightarrow\sigma^*V$ constructed from
the coefficient blocks equals $p^3\Phi^{-1}$.
\end{lemma}
\begin{proof}
The inverse of $C$ is a Frobenius intertwiner on a strict
neighborhood, by the determinant formula and the construction above.
Kedlaya's uniqueness theorem for hypergeometric intertwiners
\cite[Lemma~2.3.6 and Remark~4.1.3]{IV:kedlaya} applies. In the
coordinate $\lambda z$ the upper parameters are $1/3,1/3,2/3,2/3$
and the lower parameters are $1,1,1,1$. No upper-minus-lower
difference is integral. The singularities are distinct modulo $p$,
the exponents are $p$-integral, and multiplication by $p$ permutes
the upper classes and fixes the lower classes modulo $\Z$. The
coordinate rescaling changes the Frobenius lift by the unit
$\lambda^{1-p}\equiv1\bmod p$, which is allowed by the Taylor
change-of-lift formula \cite[Remark~2.3.7]{IV:kedlaya}. Hence
$C=\kappa\,p^3\Phi^{-1}$ for some $\kappa\in\Q_p^\times$.

At any smooth rational Teichm\"uller parameter $\widehat z$, where
$\widehat z^p=\widehat z$, the determinant of $C$ is $p^6$.
The alternating perfect pairing of the polarized rank-four $V$ has
similitude $p^3$ so $\det\Phi=p^6$. Thus $\kappa^4=1$.
To determine which root it is, let
\[
 P(z)=\sum_{0\le k\le\lfloor(p-1)/3\rfloor} A_k^2 z^k\quad\text{in }\F_p[z].
\]
At a Teichm\"uller parameter the row estimates of the explicit matrix
give $\Tr C\equiv P(z)\pmod p$. The same holds for $\Tr\Phi$.
Here is the point-count justification over all smooth parameters.
In the first Hesse pencil use the three coordinates proportional to
$(X^3,Y^3,Z^3)$ and summing to one; in the second use the coordinates
proportional to $(X^3,Y^3,z Z^3)$ and summing to one. Their respective
products are $u$ and $z/u$. The degree-three quotient in each
elliptic fiber and the same cyclic descent as in
Appendix~\ref{IV:app:torus} identify the rank-four parabolic term
throughout this locus. This description is defined over the base
before choosing a cube root of $z$.
The product-moment count there, with $z$ in place of $c$, gives
$\Tr\Phi\equiv P(z)$. Self-duality pairs every eigenvalue
$\alpha$ with $p^3/\alpha$, so $\Tr(p^3\Phi^{-1})=\Tr\Phi$.
Finally $P(0)=1$ and $\deg P\le(p-1)/3<p-2$; at least one of the
$p-2$ smooth nonzero residues has $P(z)\ne0$. Reduction there gives
$\kappa\equiv1\bmod p$. As $p\nmid4$, this forces $\kappa=1$.
\end{proof}

\section{The rational lift and the third digit}\label{IV:sec:fiber}
We require the value at the characteristic-zero parameter $c$, whose
de Rham holomorphic vector lies in the geometric rank-two space $W$.
We therefore compare the fibers at $c^p$ and $c$. Put
$L=\log(c^{p-1})\in p\Z_p$. In the cyclic coordinates at $c$, Taylor
transport from $c^p$ to $c$ sends the vector $\theta^j\Omega$ evaluated
at $c^p$ to
\[
 \sum_{h\ge0}\frac{L^h}{h!}\,(\theta^{h+j}\Omega)(c).
\]
This follows by pairing against a horizontal section in the residue
disk and expanding its scalar functions at $c$ in the coordinate
$\log(z/c)$. The same formula follows directly by Taylor transport
for the Gauss--Manin connection. The companion relation reduces
every $\theta^n\Omega$ integrally in the cyclic lattice at $c$, since
$1-\lambda c=9/8$ is a $p$-adic unit. The series converges, and
$v_p(L^h/h!)\geq\min(h,3)$ for $p\ge5$. Denote the transported
endomorphism of $V_c$ by $C_c$.

\begin{lemma}\label{IV:lem:div}
In the ordered basis $(e_0,e_1,e_2,e_3)$ at $c$, with $e_i=\theta^i\Omega$,
put $\Lambda=\sum_i\Z_p e_i$. Then
\[
 C_c=\operatorname{diag}(1,p,p^2,p^3)\,U,\qquad
 U\in\operatorname{GL}_4(\Z_p),\qquad
 (C_c)_{00}\equiv S_p\pmod{p^3}.
\]
\end{lemma}
\begin{proof}
Before transport, the coefficient of $\theta^j\Omega$ in the
unreduced series of each column is divisible by $p^{\min(j,3)}$.
Multiplication by the Taylor series preserves this estimate: at total
order $j+h$ the valuations add to at least $\min(j+h,3)$.
All terms of order three or more are integral multiples of $p^3$
after reduction by the companion equation. Orders zero, one and two
are already basis vectors. This proves the four row divisibilities.
For the first column the order-zero coefficient is exactly
$\sum_{k<p}c^k b_k(0)=S_p$. Lower-order Taylor corrections involve
$e_1$ and $e_2$ only. Thus reduction of the first coordinate changes
$S_p$ by a multiple of $p^3$. By Lemma~\ref{IV:lem:normal}, $C_c$ is the
geometric $p^3\Phi_c^{-1}$, and its determinant is $p^6$ at this
fiber as well. Dividing row $i$ by $p^i$ gives an integral matrix
of determinant one.
\end{proof}

\begin{proof}[Proof of Theorem~\ref{IV:thm:main}]
The averaging correspondence in Appendix~\ref{IV:app:tetra} is defined
over $\Q$ on $V_c$ and has good reduction for $p\ge5$. Functorial
crystalline comparison at the lift $c$ identifies its de Rham image
with a Frobenius-stable rank-two space. Its holomorphic line contains
$e_0$. The comparison in Appendix~\ref{IV:app:modular} gives the
characteristic polynomial of $\Phi_c$ on this space as
$X^2-a_pX+p^3$. The same polynomial holds for $C_c=p^3\Phi_c^{-1}$,
since the two roots have product $p^3$. Cayley--Hamilton applied on
this space gives the vector identity
\[
        C_c e_0+p^3 C_c^{-1}e_0=a_p e_0 .
\]
By Lemma~\ref{IV:lem:div}, $C_c^{-1}e_0=U^{-1}e_0\in\Lambda$.
It follows that $C_c e_0\equiv a_p e_0\bmod p^3\Lambda$.
Taking its first coordinate and using the same lemma proves the
congruence. This argument works for all good primes, independently
of the slopes of this rank-two factor: the integrality used here
is in the explicit cyclic-coordinate lattice.
\end{proof}

\begin{corollary}\label{IV:cor:harmonic}
Put $m=\lfloor(p-1)/3\rfloor$ and $q_p(c)=(c^{p-1}-1)/p$.
With $H_j=\sum_{i=1}^{j}1/i$ and $H_0=0$, every good prime satisfies
\[
 \sum_{k=0}^{m} A_k^2 c^k
       \bigl(q_p(c)+6(H_{3k}-H_k)\bigr)\equiv0\pmod p .
\]
\end{corollary}
\begin{proof}
Read row one of $C_c e_0\equiv a_p e_0\bmod p^3\Lambda$ at precision
$p^2$. Its unreduced coefficient is the coefficient of $T$ in
$\exp(LT)\sum_{k<p}c^k b_k(T)$. Reducing higher powers of $\theta$
does not affect this coefficient modulo $p^2$. We have
$L\equiv p q_p(c)\bmod p^2$ and $b_0'(0)=0$. For $3k<p$,
logarithmic differentiation of the finite factorial quotient gives
$b_k'(0)=6pA_k^2(H_{3k}-H_k)$. For $3k\ge p$ at least one squared
factor $(p h+3pT)^2$ occurs in the numerator and the entire series
is divisible by $p^2$. Dividing the row-one congruence by $p$
proves the claim without any division by $S_p$.
\end{proof}
\section{Exact checks}\label{IV:sec:checks}
The accompanying code tests the Frobenius matrix and the newform independently.
The first tests the transported matrix at every smooth nonzero
parameter in $\F_p$ for the 23 primes $5\le p\le97$ (1,009 fibers).
It compares its trace modulo $p^3$ with an independent count of pairs
of toric cubics. At $z=c$ it also tests the vanishing of its two
off-diagonal first-column entries modulo $p^3$ in rows one and two;
row three is divisible by $p^3$ by construction.

The second computes the Fourier coefficients of $f$ without using
the hypergeometric period. The 21 eta quotients used in
\texttt{verify\_finite.py} give a basis for $S_4(\Gamma_0(54))$;
the checks of their cusp orders, character and independence, and the
Sturm bound 36, are made over $\Q$. Imposing $U_2=2$, $U_3=0$ and
$T_5=3$ leaves a one-dimensional joint eigenspace, which is normalized
by its leading coefficient. For the longer computation we expand twice
this form modulo $2^{64}$ using Euler's pentagonal number theorem.
The bound $|a_p|\le2p^{3/2}$ makes the signed lift unique in the
tested range. The ratio of successive truncation terms is
$-(3k-1)^2(3k-2)^2/(648k^4)$; clearing prefix denominators gives a
division-free check modulo $p^3$. Terms with $3k\ge2p$ vanish at
this precision.
All 9,590 primes $5\le p\le100{,}000$ passed. For example,
$a_{99991}=49{,}718{,}099$. In the separate computation modulo $p^4$,
all 667 primes $5\le p\le5000$ had $v_p(S_p-a_p)=3$.
Already $S_5\equiv128\bmod625$ and $a_5=3$ prove that the exponent
three cannot be increased uniformly. Every prime up to $100{,}000$
in this computation had $p\nmid a_p$; the proof of the theorem makes
no assumption about the existence or distribution of nonordinary primes.

\appendix\paperanchors{IV.app}
\section{The tetrahedral projector}\label{IV:app:tetra}
\subsection{The Hesse cover}
We also write $z_*=c$.
B\"onisch~\cite[\S5.1]{IV:bonisch} constructs an elliptic surface $S_z$ with
equation $s_0XYZ=t_0(X^3+Y^3+zZ^3)$ and forms
$S_1\times_{\PP^1}\epsilon^*S_z$, where $\epsilon$ interchanges
$s_0$ and $t_0$. Set $s=s_0/(3t_0)$. Since $c=(-1/18)^3$,
putting $Z_{\rm new}=(-1/18)Z_{\rm old}$ in the second cubic
identifies this threefold over $\Q$ on the smooth locus with the product of
\[
 H_s:\ X^3+Y^3+Z^3=3sXYZ,\qquad H_{I(s)},\qquad I(s)=-2/s .
\]
Let $D=\{\infty,1,\zeta,\zeta^2,0,-2,-2\zeta,-2\zeta^2\}$, where
$\zeta$ is a primitive third root of unity. On
$j:\PP^1\setminus D\hookrightarrow\PP^1$ write $L=R^1\pi_*\Q_\ell$
for the first Hesse local system, and set
\begin{equation}\label{IV:eq:par}
 U=H^1(\PP^1_{\overline{\Q}},j_*(L\otimes I^* L)).
\end{equation}
This is the nonvertical component $U_{z_*}$ in
\cite[\S\S4.3, 5.1]{IV:bonisch}. Here are the details of that
identification. In Leray for the threefold over $\PP^1$ the term
$E_2^{1,2}$ has no incoming or outgoing differentials.
Fiberwise K\"unneth expresses $R^2$ as the two elliptic $R^2$
sheaves and $j_*(L\otimes I^*L)$, since at most one factor is
singular over each point. Each elliptic $R^2$ maps onto the
constant Tate line, with skyscraper kernel, and so has zero $H^1$.
Furthermore $E_2^{2,1}=0$: here $R^1=j_*(L\oplus I^*L)$, and the
dual elliptic systems have no global invariants. The Leray filtration
therefore identifies $E_2^{1,2}$ with a subspace of $H^3$ annihilated
by restriction to fibers. It lies in B\"onisch's eight-dimensional
kernel of the vertical component restrictions. Equality follows by
the dimension calculation for parabolic cohomology below. The
Euler characteristic gives $\dim U=8$: the tensor has rank
four, its invariant stalk at each puncture has dimension two, and its
global invariants and dual invariants vanish. Indeed, the two
irreducible elliptic local systems have different bad loci, hence are
not geometrically isomorphic, even after dualizing. This also gives
$H^0(\PP^1,j_*(L\otimes I^*L))=H^2(\PP^1,j_*(L\otimes I^*L))=0$.

\subsection{Two simultaneous transformations}
Consider
\[
 d(s)=\zeta s,\qquad m(s)=\frac{s+2}{s-1}.
\]
Then $I d=d^{-1} I$ and $I m=m I$. On projective triples let
$D_0=\operatorname{diag}(\zeta^2,1,1)$ and
$F=(\zeta^{ij})_{0\leq i,j\leq2}$.
Lifts on the pair of elliptic curves are $(D_0,\bar D_0)$ for $d$ and
$(F,\bar F)$ for $m$. For the latter this follows by expanding the cube
sum and product of the three entries of $F(X,Y,Z)^{\mathsf t}$:
they are $3(X^3+Y^3+Z^3+6XYZ)$ and $X^3+Y^3+Z^3-3XYZ$.
Thus the new Hesse parameter is $m(s)$.
These projectivities preserve the base locus of the pencil as a set, so
also act on the resolved rational elliptic surfaces.
On $\{\infty,1,\zeta,\zeta^2\}$ the two base maps act as a three-cycle
and a double transposition. They generate $A_4$.

There is a minor distinction between matrices on the pencil and their
actions on the tensor. A word acting trivially on the base induces an
automorphism of the generic elliptic curve, which is a translation
followed by $[1]$ or $[-1]$, since its $j$-invariant is nonconstant.
The sign on the second factor is the same: the second word is obtained
from the first by conjugating $\zeta$ and evaluating at $I(s)$, and
both signs are rational. Thus such a word acts trivially on
$L\otimes I^*L$. We obtain an action of $A_4$ on $U$ and on its
Betti and de Rham counterparts. Conjugation sends $d$ to $d^{-1}$.
It also preserves the action of $m$: conjugation interchanges $F$ and
$\bar F=FP$, where $P$ exchanges the last two coordinates. The
involution $P$ acts by $-1$ on $H^1$; it fixes a point and is
nontrivial on the generic fiber. The two signs cancel on the tensor.
Consequently the two projectors
\begin{equation}\label{IV:eq:avg}
 e_C=\frac{1+d+d^2}{3},\qquad
 e=\frac1{12}\sum_{h\in A_4}h
\end{equation}
commute with $G_\Q$. The averaging takes place on cohomology by
algebraic correspondences; $e_Ce=ee_C=e$. B\"onisch identifies
$V=e_CU$ with $V_{z_*}$, of Hodge numbers $(1,1,1,1)$.

\subsection{The two ranks and their Hodge types}
B\"onisch computes $\Tr(d\mid U)=\Tr(d^{-1}\mid U)=2$.
For $m$ the only fixed base points are $1\pm\sqrt3$, both smooth.
At either point $F^2=3P$ in homogeneous coordinates, and hence the
square of the fiber action on $H^1$ is $-1$. The action has determinant
$1$, since it is an elliptic automorphism of degree one. Its eigenvalues
are $i,-i$, and its trace vanishes. The conjugate factor has trace
zero as well. The two fixed points are simple. The Lefschetz formula
for the equivariant constructible sheaf in~\eqref{IV:eq:par} now yields
$\Tr(m\mid U)=0$, using the vanishing of $H^0$ and $H^2$.
These are character traces for finite-order geometric symmetries.
The identity, the three involutions and the two classes of four
elements of order three give
\[
 \dim eU=\frac{8+3\cdot0+4\cdot2+4\cdot2}{12}=2.
\]
Write $W=eU$ and $K=(e_C-e)U$. The latter also has dimension two.
The line of type $(3,0)$ in $U$ is fixed pointwise by $d$ by
\cite[\S5.1]{IV:bonisch}. The entire $A_4$ preserves this Hodge line, so
its action there factors through $A_4^{\rm ab}\simeq C_3$, generated
by the class of $d$. Thus $W$ has types $(3,0),(0,3)$, and $K$ has
types $(2,1),(1,2)$. The group acts by isometries of the natural
alternating pairing (cup product on parabolic cohomology). The
averages are self-adjoint, so the restrictions to $W$ and $K$ are
nondegenerate and both determinants are $\Q_\ell(-3)$.
\section{Frobenius and the torus count}\label{IV:app:torus}
For $p\geq5$ the reduction is good. In Hesse coordinates the two
sets of four singular fibers remain disjoint: the only possible finite
overlap between $\mu_3$ and $-2\mu_3$ would imply $1=-8$.
At these primes the projectors~\eqref{IV:eq:avg} have $p$-integral
coefficients. Their rational crystalline realizations are weakly
admissible and compatible with the good-prime traces on the
$\ell$-adic realization. In particular
\begin{equation}\label{IV:eq:slopes}
 v_p(\Tr(\Frob_p\mid K))\geq1 .
\end{equation}
Indeed $K$ has Hodge slopes 1 and 2. The Newton polygon lies above
the Hodge polygon, so each Frobenius eigenvalue has valuation at least
one. Their product equals $p^3$. This observation also distinguishes
the factors whenever only one of the two candidate traces is divisible
by $p$.


\subsection{Trace on the four-dimensional constituent}
On the open Hesse pencil define
\[
 (x,y,z)=\frac{(X^3,Y^3,Z^3)}{3sXYZ}.
\]
Then $x+y+z=1$, $xyz=u=1/(27s^3)$. On the paired pencil the
corresponding parameter is $v=c/u$. This is a degree-three isogeny
after a choice of origin (equivalently, a finite map between the
smooth genus-one compactifications inducing an isomorphism on
rational first cohomology). Indeed the three preimages form an orbit
of the order-three diagonal projectivity
$(X,Y,Z)\mapsto (X,\zeta Y,\zeta^2Z)$ on the generic fiber.
The target coordinates are unchanged by $s\mapsto\zeta s$ with its
specified lift $D_0$, and the paired coordinates by its conjugate.
Taking the invariants of this cyclic base cover consequently identifies
$V$ with the parabolic cohomology in the $u$ coordinate of the
product of the two elliptic local systems. One may check the descent
on the open locus and then extend by invariant stalks. Finite
pushforward and projection onto invariants commute with this
extension in characteristic zero.

Write $N(u)$ for the number of triples in $(\F_p^\times)^3$ whose
sum is $1$ and product is $u$. The cubic compactification
$x y(w-x-y)=u w^3$ adds three rational points at infinity. Its smooth
trace and its nodal invariant-stalk trace at $u=1/27$ therefore both
equal $t(u)=p-2-N(u)$ for $u\in\F_p^\times$. For the nodal cubic
this follows either by normalization or by counting its smooth
projective model with the two branches identified.
There is no invariant stalk of the product at $0$ or $\infty$.
To see this, a Weierstrass model for the first elliptic system is
$Y^2+XY+uY=X^3$, obtained with $X=-u/x$ and $Y=-Xy$. Its invariants
are $c_4=1-24u$ and $\Delta=u^3(1-27u)$. Thus at $u=0$ its
monodromy is unipotent. At infinity the substitutions $h=1/u$,
$X=h^{-2}X'$, $Y=h^{-3}Y'$ give minimal discriminant of order eight
and $c_4$ of order three (Kodaira type $\mathrm{IV}^*$). Its tame
inertia on first cohomology has eigenvalues $\zeta,\zeta^2$. In the
paired system these roles are interchanged at the two points.
At the two distinct nonzero nodes only one factor is singular.
The trace formula, with no global invariants or dual invariants, gives
\begin{align}
 T_p:=\Tr(\Frob_p\mid V)
   &=-\sum_{u\in\F_p^\times}t(u)t(c/u) \notag\\
   &=p^3-5p^2+10p-8-M_p,\label{IV:eq:exacttrace}\\
 M_p&=\sum_{u\in\F_p^\times}N(u)N(c/u).\notag
\end{align}
In obtaining~\eqref{IV:eq:exacttrace} we have used
$\sum_uN(u)=p^2-3p+3$, by inclusion and exclusion for the three
affine lines $x=0$, $y=0$ and $x+y=1$.

\subsection{Extracting the truncated period}
Here $M_p$ counts ordered pairs of triples. On six nonzero
coordinates with the prescribed two sums, write $u_0$ for their total
product. Modulo $p$ the indicator of $u_0=c$ is
$-\sum_{j=1}^{p-1}(u_0/c)^j$. Put
$B_j=\sum_{x+y+z=1,\,xyz\ne0}(xyz)^j$. Then
\[
 M_p\equiv-\sum_{j=1}^{p-1}c^{-j}B_j^2\pmod p.
\]
For $j=p-1$ we have $B_j\equiv3$. For $1\leq j\leq p-2$ put
$k=p-1-j$. Expanding $x^jy^j(1-x-y)^j$ and summing over all $x,y$
leaves only the monomial with added exponents $k,k$. It occurs exactly
when $3k\leq p-1$, and has coefficient
\[
 \frac{(p-1-k)!}{(k!)^2(p-1-3k)!}\equiv A_k\pmod p .
\]
The two nonzero power sums each contribute $-1$. For the remaining
$k$ the coefficient $A_k$ is zero modulo $p$. It follows that
\begin{equation}\label{IV:eq:modp}
 M_p\equiv -8-S_p\pmod p,\qquad T_p\equiv S_p\pmod p .
\end{equation}
In particular, with the identification of $W$ with $f$, equations
\eqref{IV:eq:slopes}--\eqref{IV:eq:modp} give $S_p\equiv a_p\pmod p$
at every good prime, including any nonordinary prime.

\section{Modular identification}\label{IV:app:modular}
Let $g$ denote the normalized newform $54.2.\mathrm{a}.\mathrm{a}$ with good-prime coefficients $b_p$.
\begin{theorem}\label{IV:thm:comparison}
For the rank-four constituent of the specified Hesse model, the
semisimplifications satisfy
\[
 V_{c,\ell}^{\mathrm{ss}}\simeq
 \rho_{f,\ell}\oplus\rho_{g,\ell}(-1).
\]
The good-prime Euler polynomial of $V$ is therefore
$(1-a_pT+p^3T^2)(1-pb_pT+p^3T^2)$.
\end{theorem}
\begin{proof}
Compare $W_2$ with $\rho_{f,2}$ and $K_2(1)$ with $\rho_{g,2}$.
Each pair has the same determinant by the alternating pairing and
the trivial nebentypus of the newforms. Their ramification is confined
to $2,3$: the two sets of bad Hesse fibers are disjoint and tame
over $\Z[1/6]$ after adjoining $\zeta$, the rational elliptic surfaces
are smooth there, and the fiber product has good reduction.
Compactness gives stable $\Z_2$-lattices even though the averaging
operators were defined over $\Q_2$.

The four polynomials can be determined by point counts over $\F_p$
and $\F_{p^2}$. Write $t_q=\Tr(\Frob_q\mid V)$. The trace-count
identity~\eqref{IV:eq:exacttrace} also applies over these extensions.
The reciprocal quartic has coefficients $1,-t_p,(t_p^2-t_{p^2})/2,
-p^3t_p,p^6$ in ascending degrees. The following exact counts are
reproducible by the accompanying \texttt{verify\_four\_primes.py}.
They agree with B\"onisch's computations in \cite[\S5.1]{IV:bonisch}.
\[
\begin{array}{r|rr|rr|rr}
p&t_p&t_{p^2}&\Tr W&\Tr K&a_p& b_p\\ \hline
5& 18&-266&3&15&3&3\\
7& 22&-482&29&-7&29&-1\\
13&-32&-5684&20&-52&20&-4\\
19&-68&-14756&-106&38&-106&2
\end{array}
\]
Here is why the indicated ordering of factors is forced.
The two displayed reciprocal pairs have distinct traces and each
has product $p^3$. A root of the quartic has just one reciprocal
partner, so its two rank-two summands have to be those pairs.
The minimum Hodge slope of $K$ is one. Crystalline comparison at
the four good primes and the Newton--Hodge inequality therefore give
$v_p(\lambda)\geq1$
for both of its roots. Precisely one of the two candidate traces is
divisible by $p$. The use of the same projected factor across
realizations follows from the geometric action: the two projectors
are rational combinations of automorphisms on the parabolic term.
More explicitly, on the smooth projective fiber product one takes
the alternating trace of an automorphism composed with Frobenius by
intersection with the diagonal. Degrees $0,2,4,6$ are algebraic and
degrees $1,5$ vanish. To obtain the primitive term in degree three
one removes the ruled-fiber contribution, given by the reduced
component spaces tensored with the elliptic $H^1(-1)$. This
calculation is common to the crystalline and $\ell$-adic realizations.
If Frobenius permutes automorphisms or components, take the trace
of the entire averaged correspondence. The last two columns are the
Fourier coefficients of the labeled newforms \cite[\S5.1]{IV:bonisch}.

We specify the finite comparison criterion, since agreement of traces
at an unspecified set of primes would not suffice. For each of the
four representations, the residual characteristic polynomials at
$5,7,13$ are $X^2+X+1,X^2+X+1,X^2+1$ respectively.
Li, Long and Yang, Lemma~6.1 in \cite{IV:lly}, show that these data and
the ramification bound determine a unique surjective residual
representation up to conjugacy. Its field is the splitting field
of $X^3+3X-2$, with quadratic subfield $\Q(\sqrt{-6})$.
Their Theorem~6.1 and its field analysis in the proof of
Theorem~6.2 now apply to each pair. In detail, a nonzero deviation
between equal-determinant lifts would be detected at a Frobenius
of order at least four in one of the deviation groups
$S_4,S_4\times C_2,S_3\times C_2$.
For the first two possibilities the three admissible quartic
fields have defining polynomials
\[
 X^4-4X-3,\qquad X^4-8X+6,\qquad X^4-12X^2-16X+12 .
\]
The first two yield a four-cycle at $13$, the third at $19$.
In the remaining case an independent quadratic character unramified
outside $2,3$ is nontrivial at $5$ or $7$: among the squareclasses
generated by $-1,2,3$ only $1,-6$ split at both primes. The residual
Frobenius there has order three, giving an element of order six.
The table has equal traces for each comparison at all four primes,
eliminating these deviations. One may extend scalars to
$\Q_2(i)$ to use precisely the coefficient ring of \cite{IV:lly};
its residue field is still $\F_2$.

This proves both comparisons at $\ell=2$, hence the asserted
Euler polynomials at every $p\geq5$. The common Frobenius
polynomials for $V$ across $\ell$, followed by Chebotarev and
Brauer--Nesbitt, give the stated semisimplifications for all $\ell$.
\end{proof}

\section{Finite product identities for the two higher digits}\label{IV:app:digits}
By Theorem~\ref{IV:thm:comparison} and \eqref{IV:eq:modp} the first digit is established for all
good primes. The congruence modulo $p^3$ is equivalent, without an
ordinarity hypothesis, to the scalar condition
\begin{equation}\label{IV:eq:two}
 \frac{M_p+8+S_p}{p}+b_p-10+5p\equiv0\pmod{p^2}.
\end{equation}
The left quotient is integral by \eqref{IV:eq:modp}. The equivalence
follows by substituting $a_p=T_p-pb_p$ in \eqref{IV:eq:exacttrace}.

\subsection{Jacobi sums and factorials}
Let $p\geq5$, $c=-1/5832$, and $T_p=\Tr(\Frob_p\mid V_{z_*})=a_p+p b_p$.
Set $m=\lfloor(p-1)/3\rfloor$, $n=\lfloor 2(p-1)/3\rfloor$, and, for $0\le k\le m$, put
\[
 \delta_k=\frac{pk}{p-1},\qquad
 W_k=\frac{\prod_{j=1}^{3k}(j-3\delta_k)}
              {\prod_{j=1}^{k}(j-\delta_k)^3}.
\]
All the denominators in $W_k$ are $p$-adic units. For $m< k\le n$ set
$r_k=3k-(p-1)$. Define the following explicitly computable $p$-adic integer:
\[
 B_p=\frac1{1-p}\left(
 \sum_{k=0}^{m} c^{p^2 k} W_k^2
 +p^2\sum_{k=m+1}^{n} c^{k}
                 \left(\frac{r_k!}{(k!)^3}\right)^{\!2}\right).
\]
\begin{proposition}\label{IV:prop:gamma-digits}
The good-prime trace satisfies $T_p\equiv B_p\pmod {p^3}$. In particular,
$B_p-S_p\in p\Z_p$, and the congruence $S_p\equiv a_p\pmod {p^3}$
for a given $p$ is equivalent to
\[
 \frac{B_p-S_p}{p}\equiv b_p\pmod {p^2}.
\]
Write $H_j=\sum_{i=1}^j i^{-1}$, $H_0=0$, and $q_p(c)=(c^{p-1}-1)/p$.
At precision $p^2$ one obtains the equivalent identity in $\F_p$
\[
 b_p=\sum_{k=0}^{m}A_k^2 c^k\left[1+k q_p(c)+6k(H_{3k}-H_k)\right].
\]
\end{proposition}
\begin{proof}
For $u\in\F_p^\times$ put
\[
 N_u=\#\{(x,y,w)\in(\F_p^\times)^3:x+y+w=1,\ xyw=u\},\quad h(u)=p-2-N_u.
\]
The toric trace formula gives $T_p=-\sum_{u\ne0}h(u)h(c/u)$.
We have $\sum_{u\ne0}N_u=p^2-3p+3$, so $\sum_{u\ne0}h(u)=-1$.
For a nontrivial multiplicative character $\chi$, its Mellin coefficient
$\sum_u h(u)\chi(u)$ is $-J(\chi,\chi,\chi)$, by the defining equation for $N_u$.
With $\omega$ the Teichm\"uller character, write
$J_k=J(\omega^{-k},\omega^{-k},\omega^{-k})$ for $1\le k\le p-2$.
Character orthogonality gives the exact identity
\[
 T_p=\frac1{1-p}\left(1+\sum_{k=1}^{p-2}\omega(c)^k J_k^2\right).
\]
For each $k\ge1$ write $3k=e(p-1)+r$ with $1\le r\le p-1$, and put
$x=k/(p-1)$. The Gross--Koblitz formula (B.~H.~Gross and N.~Koblitz,
\emph{Gauss sums and the $p$-adic $\Gamma$-function}, Ann. of Math.
109 (1979), 569--581, Theorem~1.7)
for the usual Gauss sum is $g(\omega^{-k})=-\pi^k\Gamma_p(x)$ with $\pi^{p-1}=-p$.
The Gauss--Jacobi sum relation gives
\[
 J_k^2=p^{2e}\left(\frac{\Gamma_p(x)^3}{\Gamma_p(r/(p-1))}\right)^{\!2}.
\]
Indeed, for $r<p-1$ the unsquared Jacobi sum is $g(\omega^{-k})^3/g(\omega^{-r})$.
For $r=p-1$ it is $-g(\omega^{-k})^3/p$, which gives the same squared expression since
$\Gamma_p(1)=-1$. This convention includes both integral break points of $3k/(p-1)$.
All terms with $e\ge2$ vanish modulo $p^3$.

For $e=0$ write $x=-k+\delta_k$ and $3x=-3k+3\delta_k$. Morita's functional equation and $3k\le p-1$ give
\[
 \frac{\Gamma_p(x)^3}{\Gamma_p(3x)}=W_k\,\frac{\Gamma_p(\delta_k)^3}{\Gamma_p(3\delta_k)}.
\]
For $v\in p\Z_p$ one has $\Gamma_p(v)^3/\Gamma_p(3v)\equiv1\pmod{p^3}$.
To see this, each block $\prod_{j=1}^{p-1}(j+hp)$ is congruent to $(p-1)!$ modulo $p^3$ for integers $h\ge0$:
$H_{p-1}\equiv0\pmod{p^2}$ by Wolstenholme, and $\sum_{j=1}^{p-1}j^{-2}\equiv0\pmod p$.
Thus $\Gamma_p(hp)\equiv\Gamma_p(p)^h\pmod{p^3}$ by Morita's product definition.
The ratio at $v=hp$ equals $1$ modulo $p^3$, and continuity gives the assertion for all $v\in p\Z_p$.
For $e=1$ the gamma ratio is needed only modulo $p$. Morita's function is constant modulo $p$ on each residue class modulo $p$ (this also follows by grouping its integral product into blocks).
Since $x\equiv-k$ and $r/(p-1)\equiv-r\pmod p$, and since
$\Gamma_p(-j)=1/j!$ for $0\le j\le p-1$, this ratio is $r!/(k!)^3$ modulo $p$.
Finally $\omega(c)\equiv c^{p^2}\pmod{p^3}$. These calculations give $T_p\equiv B_p\pmod{p^3}$.

For the second digit, $\delta_k\equiv-pk\pmod {p^2}$, and expansion of the finite products yields
\[
 W_k^2\equiv A_k^2\bigl[1+6pk(H_{3k}-H_k)\bigr]\pmod {p^2}.
\]
Also $c^{p^2 k}\equiv c^k(1+pk q_p(c))\pmod{p^2}$. The terms of $S_p$ with $k>m$ are in $p^2\Z_p$.
Expanding $(1-p)^{-1}$ now gives the displayed identity for the obstruction modulo $p$. The two equivalences follow by substituting $T_p=a_p+p b_p$.
\end{proof}

\subsection{The elliptic trace as a second member of the same pencil}
The curve with LMFDB label $54.\mathrm{a}3$ corresponds to $g$ and has
minimal equation $y^2+xy=x^3-x^2+12x+8$. Its rational point $(1,4)$
has order three \cite{IV:curve}. The substitutions $x=1+X$,
$y=4+X+Y$ give
\[
 Y^2+3XY+9Y=X^3 .
\]
On replacing $(X,Y)$ by $(9X',27Y')$, one obtains
$Y'^2+X'Y'+\tfrac13Y'=X'^3$. Thus the weight-two trace is precisely
the trace $t(1/3)$ of the toric cubic in Appendix~\ref{IV:app:torus}. The
coordinate changes have good reduction for $p\geq5$.

In fact for each $u\in\F_p^\times$ Fourier inversion for a single
copy of $N(u)$, with the coefficients $B_j$ above, gives
\[
 N(u)\equiv -3-\sum_{k=1}^m A_k u^k,\qquad
 t(u)\equiv \sum_{k=0}^m A_k u^k\pmod p .
\]
Write $L_p=\sum_{k=0}^m A_k3^{-k}$ in $\F_p$.
The second-digit identity is therefore the purely finite congruence
\begin{equation}\label{IV:eq:pure}
 \sum_{k=0}^m A_k^2c^k\bigl[1+k q_p(c)+6k(H_{3k}-H_k)\bigr]
 = L_p\quad \text{in }\F_p .
\end{equation}
This includes all good primes, with no hypothesis on $a_p$.
For the third digit one can even recover the integer $b_p$ from $L_p$:
it is the unique integer $b$ satisfying
\[
 b\equiv L_p\pmod p,\qquad b\equiv p+1\pmod3,\qquad |b|\leq2\sqrt p .
\]
Indeed the rational point of order three injects under good
reduction, giving the second residue condition. The Hasse bound
gives the interval, and $3p>4\sqrt p$ ensures uniqueness.
Consequently Proposition~\ref{IV:prop:gamma-digits} gives an all-prime
criterion using only factorials and elementary arithmetic, independent
of any tables of newforms. There is also a version which involves no reconstruction of an integer.
Put $v=1/3$ and $R_k=r_k!/(k!)^3$ on the middle range.
By inversion for the single trace $h(v)$ the Mellin formula reads
\[
 b_p=\frac{1}{1-p}\left(1+\sum_{k=1}^{p-2}\omega(v)^k J_k\right).
\]
In the notation of the preceding proof the unsquared Jacobi sum
equals $(-p)^e\Gamma_p(x)^3/\Gamma_p(r/(p-1))$, including $r=p-1$.
Reducing to modulus $p^2$ yields
\[
 (1-p)b_p\equiv \sum_{k=0}^{m}v^{pk}W_k
                     -p\sum_{k=m+1}^{n}v^kR_k\pmod{p^2}.
\]
Together with Proposition~\ref{IV:prop:gamma-digits}, this proves
that the good-prime conjecture is equivalent to the following single
finite product identity for each $p\geq5$:
\begin{equation}\label{IV:eq:fullpure}
 \sum_{k=0}^{m}\big(c^{p^2k}W_k^2-pv^{pk}W_k\big)
 +p^2\sum_{k=m+1}^{n}\big(c^kR_k^2+v^kR_k\big)
 \equiv(1-p)S_p\pmod{p^3}.
\end{equation}
By Theorem~\ref{IV:thm:main} both \eqref{IV:eq:pure} and \eqref{IV:eq:fullpure} hold for every prime $p\ge5$.


\subsection{Computational details for the finite identities}
The index of $\Gamma_0(54)$ is 108, with 12 cusps, genus four and
no elliptic points. Thus $\dim S_4(\Gamma_0(54))=21$ and the Sturm
bound is 36. In the basis of eta quotients used in Section~\ref{IV:sec:checks},
imposing $U_2=2$, $U_3=0$ and $T_5=3$ gives subspaces of dimensions
$4,3,1$, respectively. At $p=4999$ the normalized eigenform gives
$a_p=-560608$ and the independent truncation gives
$S_p\equiv124924454391\pmod{p^3}$.
The Jacobi-sum and elliptic-trace calculations were checked separately
by point counts at all 93 primes $5\le p\le500$.


\endgroup
\renewcommand{\runningpaper}{References}
\phantomsection\addcontentsline{toc}{part}{References}
\begin{thebibliography}{99}
\bibitem{I:BD}
F.~Beukers and E.~Delaygue,
\emph{Some supercongruences of arbitrary length},
Indagationes Mathematicae \textbf{33} (2022), no.~5, 946--955.
\href{https://doi.org/10.1016/j.indag.2022.04.002}{doi:10.1016/j.indag.2022.04.002}.
Preprint: \href{https://arxiv.org/abs/1805.02467}{arXiv:1805.02467}. See Conjectures~1.2 and~1.5.

\bibitem{I:BM}
P.~Berthelot and W.~Messing,
\emph{Th\'eorie de Dieudonn\'e cristalline I},
Ast\'erisque \textbf{63} (1979), 17--37.
\href{https://numdam.org/item/AST_1979__63__17_0/}{Numdam full text}.

\bibitem{I:Katz}
N.~M.~Katz, \emph{$p$-adic properties of modular schemes and modular forms},
in Modular Functions of One Variable III,
Lecture Notes in Mathematics \textbf{350}, Springer, 1973, pp.~69--190.
\href{https://web.math.princeton.edu/~nmk/old/padicpropMFMS.pdf}{Author's full text}.

\bibitem{I:Stacks}
The Stacks Project Authors, \emph{The Stacks Project},
Proposition~60.17.4 (Tag~07JH), Taylor comparison for crystals.
\url{https://stacks.math.columbia.edu/tag/07JH}.

\bibitem{I:SunYe}
Z.-H.~Sun and D.~Ye, \emph{Supercongruences via Beukers' method},
\href{https://arxiv.org/abs/2408.09776v4}{arXiv:2408.09776v4},
19 September 2024. See Theorem~1.29 and Section~9.

\bibitem{I:BlochVanhove}
S.~Bloch and P.~Vanhove, \emph{The elliptic dilogarithm for the sunset graph},
Journal of Number Theory \textbf{148} (2015), 328--364.
\href{https://arxiv.org/abs/1309.5865}{arXiv:1309.5865}.

\bibitem{III:Beukers2026}
F. Beukers,
\emph{Supercongruences using modular forms},
Algebra Number Theory 20 (2026), 1727--1768;
\url{https://arxiv.org/abs/2403.03301}. Numbering follows arXiv v3 (16 June 2025). See Corollary~1.20, Conjecture~1.28 and Appendix Sections~8.10 and~9.5.

\bibitem{III:BV3}
F. Beukers and M. Vlasenko,
\emph{Dwork crystals III: from excellent Frobenius lifts towards supercongruences},
Int. Math. Res. Not. IMRN (2023), 20433--20483;
\url{https://arxiv.org/abs/2105.14841}.

\bibitem{III:CompleteEdition}
QEEDDD Verifiable-Mathematical-Reasoning-Corpus Project,
\emph{Excellent Frobenius Lifts and Supercongruences: Complete Collected Edition},
research manuscript, 2026.

\bibitem{III:BKSZ}
K. B\"onisch, A. Klemm, E. Scheidegger and D. Zagier,
\emph{D-brane masses at special fibres of hypergeometric families of Calabi--Yau threefolds,
modular forms, and periods},
Commun. Math. Phys. (2024);
\url{https://people.mpim-bonn.mpg.de/zagier/files/preprints/DbranesCYandMFs.pdf}.

\bibitem{III:CandelasAttractor}
P. Candelas, X. de la Ossa, M. Elmi and D. van Straten,
\emph{A one parameter family of Calabi--Yau manifolds with attractor points of rank two},
JHEP 10 (2020), 202;
\url{https://arxiv.org/abs/1912.06146}.

\bibitem{III:Fu}
Y. Fu,
\emph{$p$-adic rigidity and supercongruences at rank two attractors},
preprint (2026), \url{https://arxiv.org/abs/2609.36045}.

\bibitem{III:Katz}
N. Katz,
\emph{Exponential Sums and Differential Equations},
Annals of Mathematics Studies 124, Princeton University Press, 1990.

\bibitem{III:Deligne}
P. Deligne,
\emph{Hodge cycles on abelian varieties},
in: Hodge Cycles, Motives, and Shimura Varieties,
Lecture Notes in Mathematics 900, Springer, 1982.

\bibitem{III:Andre}
Y. Andr\'e,
\emph{Pour une th\'eorie inconditionnelle des motifs},
Publ. Math. IH\'ES 83 (1996), 5--49.

\bibitem{III:Dummigan}
N. Dummigan,
\emph{Modularity of a certain rank-2 attractor Calabi--Yau threefold},
preprint, revised 2026.

\bibitem{III:LMFDB}
The LMFDB Collaboration,
\emph{The L-functions and Modular Forms Database},
\url{https://www.lmfdb.org/}.

\bibitem{IV:bonisch}
K. B\"onisch, \emph{Modularity of special motives of
rank four associated with Calabi--Yau threefolds}, doctoral
dissertation, University of Bonn (2023), \S\S4.3 and 5.1.
\url{https://d-nb.info/1375378988/34}.

\bibitem{IV:lly}
W.-C. W. Li, L. Long and Z. Yang, \emph{On
Atkin--Swinnerton-Dyer congruence relations}, J. Number Theory
113 (2005), 117--148. See Lemma~6.1, Theorem~6.1 and the comparison
immediately preceding Theorem~6.2. \url{https://arxiv.org/abs/math/0311287}.

\bibitem{IV:kedlaya}
K. S. Kedlaya, \emph{Frobenius structures on
hypergeometric equations}, arXiv:1912.13073. See Lemma~2.3.6,
Remark~2.3.7, Proposition~3.1.11 and Remark~4.1.3.
\url{https://arxiv.org/abs/1912.13073}.

\bibitem{IV:lazda}
C. Lazda, \emph{Incarnations of Berthelot's conjecture},
J. Number Theory (2017), Corollary~6.4.
\url{https://www.math.unipd.it/~lazda/resources/Papers/Berthelot.pdf}.

\bibitem{IV:curve}
The LMFDB Collaboration, elliptic curve $54.\mathrm{a}3$,
Cremona label $54\mathrm{a}1$.
\url{https://www.lmfdb.org/EllipticCurve/Q/54/a/3}.

\bibitem{IV:lmfdb}
The LMFDB Collaboration, modular forms
\texttt{54.4.a.c} and \texttt{54.2.a.a}.
\url{https://www.lmfdb.org/ModularForm/GL2/Q/holomorphic/54/}.
\end{thebibliography}
\end{document}
